\documentclass[11pt,a4paper]{article}
\usepackage[utf8]{inputenc}
\usepackage[T1]{fontenc}
\usepackage[english]{babel}
\usepackage[margin=2.2cm]{geometry}
\usepackage{amsmath,amssymb,amsthm}
\usepackage{parskip}
\usepackage{enumitem}
\setlist{nosep,leftmargin=1.4em}
\usepackage{booktabs}
\usepackage{array}
\usepackage{xcolor}
\definecolor{linkblue}{RGB}{20,60,150}
\definecolor{defbg}{RGB}{240,244,252}
\definecolor{defframe}{RGB}{100,130,190}
\definecolor{markcol}{RGB}{176,84,0}
\definecolor{markbg}{RGB}{253,246,236}
\usepackage[most]{tcolorbox}
\newtcolorbox{defbox}[1][]{breakable,colback=defbg,colframe=defframe,boxrule=0.6pt,arc=1.5mm,left=2mm,right=2mm,top=1mm,bottom=1mm,title={#1},fonttitle=\bfseries}
\newtcolorbox{poznbox}[1][]{breakable,colback=markbg,colframe=markcol,boxrule=0.6pt,arc=1.5mm,left=2mm,right=2mm,top=1mm,bottom=1mm,title={#1},fonttitle=\bfseries}
\newtheorem{theorem}{Theorem}[section]
\theoremstyle{definition}
\newtheorem{definition}{Definition}[section]
\usepackage[colorlinks=true,linkcolor=linkblue,urlcolor=linkblue]{hyperref}
\usepackage{algorithm}
\usepackage{algpseudocode}

\floatname{algorithm}{Algorithm}
\renewcommand{\algorithmicrequire}{\textbf{Input:}}
\renewcommand{\algorithmicensure}{\textbf{Output:}}

% Quotation macro carried over from the Czech original (babel-czech \uv).
\providecommand{\uv}[1]{``#1''}

% Marker for material that is not in the EVA I / EVA II slides.
\newcommand{\mimo}{\textcolor{markcol}{\textsf{\footnotesize[$\blacklozenge$~beyond the slides]}}}
\newcommand{\mimoq}[1]{\textcolor{markcol}{\textsf{\footnotesize[$\blacklozenge$~beyond the slides: #1]}}}

\title{\textbf{Nature-Inspired Computing}\\[2mm]
\large Study Text for the Final State Examination}
\author{Final State Examination in Artificial Intelligence, MFF UK}
\date{August 2026}

\begin{document}
\maketitle

\begin{center}
\begin{minipage}{0.92\textwidth}
\small
This text covers the examination topic \emph{Nature-Inspired Computing} of the master's final
state examination in Artificial Intelligence at MFF UK. Its chapters correspond
\textbf{one to one} to the sentences of the official wording of the topic (see the mapping
table below). It is based on the slides of the lecture courses NAIL025 (Evolutionary
Algorithms~I) and NAIL086 (Evolutionary Algorithms~II), supplemented with the standard
material of the textbooks Eiben \& Smith: \emph{Introduction to Evolutionary Computing},
Mitchell: \emph{An Introduction to Genetic Algorithms}, Michalewicz: \emph{Genetic Algorithms
$+$ Data Structures $=$ Evolution Programs} and Poli, Langdon \& McPhee: \emph{A Field Guide
to Genetic Programming}. It contains definitions, pseudocode, derivations, numerical examples
and comparison tables; every chapter ends with a summary for the oral examination.
This is an English translation of the Czech study text \emph{Přírodou inspirované počítání}.
\end{minipage}
\end{center}

\begin{center}
\begin{minipage}{0.92\textwidth}
\begin{poznbox}[Scope and marking]
\small
The length of the chapters reflects the weight of the material: most space goes to genetic
algorithms with genetic programming, to schema theory, and to the applications --- that is,
to what the lectures treat in most detail.

The marker \mimo{} flags material that is \textbf{in none of the available slide decks} ---
neither EVA~I, nor EVA~II, nor the optional \emph{Evolutionary Robotics} (Mráz) --- but has
been kept in the text, either because \emph{the official wording of the topic names it}, or
because the surrounding exposition would not make sense without it. It is an offer for
deletion: whatever is not marked is backed by the slides. Where the material is only
mentioned in passing, the marker says how briefly.

The most conspicuous case is \textbf{open-ended evolution}: the topic names it explicitly,
the slides give it a single sentence. Similarly ant colony algorithms (ACO) --- the topic
speaks of swarm algorithms in the plural, the slides cover only PSO (the swarm \emph{robotics}
in Evolutionary Robotics is a different thing from swarm optimisation). Conversely,
\emph{grammatical evolution} and \emph{interactive evolution} are \emph{not} marked: they are
absent from EVA, but Evolutionary Robotics covers them.
\end{poznbox}
\end{minipage}
\end{center}

\vspace{2mm}

\section*{Mapping of the examination topic to chapters}
\addcontentsline{toc}{section}{Mapping of the examination topic to chapters}

\begin{center}
\small
\begin{tabular}{@{}>{\raggedright\arraybackslash}p{8.6cm}>{\raggedright\arraybackslash}p{7.4cm}@{}}
\toprule
\textbf{Sentence of the official topic} & \textbf{Chapter in this text} \\
\midrule
Genetic algorithms, genetic and evolutionary programming.
  & \ref{sec:ga}~Genetic algorithms, genetic and evolutionary programming \\
Schema theory, probabilistic models of the simple genetic algorithm.
  & \ref{sec:schemata}~Schema theory and probabilistic models of the SGA \\
Evolution strategies, differential evolution, coevolution, open-ended evolution.
  & \ref{sec:es}~Evolution strategies, differential evolution, coevolution, open-ended
    evolution \\
Swarm optimisation algorithms.
  & \ref{sec:roje}~Swarm optimisation algorithms \\
Memetic algorithms, hill climbing, simulated annealing.
  & \ref{sec:lokalni}~Local search and memetic algorithms \\
Applications of evolutionary algorithms (evolution of expert systems, neuroevolution, solving
combinatorial problems, multi-objective optimisation).
  & \ref{sec:aplikace}~Applications of evolutionary algorithms \\
\bottomrule
\end{tabular}
\end{center}

\vspace{2mm}
\tableofcontents
\newpage

%=====================================================================
\section{Genetic algorithms, genetic and evolutionary programming}
\label{sec:ga}
%=====================================================================

The longest sentence of the topic; the chapter therefore starts with the general scheme of
an evolutionary algorithm (without which the rest of the text makes no sense), continues
with Holland's simple genetic algorithm and the representations of an individual, and ends
with genetic and evolutionary programming, that is, the evolution of executable structures.

\subsection{The general scheme of an evolutionary algorithm}
\label{ssec:obecne}
\subsubsection{Placing the field}

\emph{Nature-inspired computing} is an umbrella term for metaheuristic algorithms
that borrow principles observed in living (and sometimes non-living) nature.
The traditional division is as follows:

\begin{itemize}
\item \textbf{Computational intelligence} (formerly \emph{soft computing}) covers
  neural networks, fuzzy systems and evolutionary computation. It stands in
  opposition to \uv{hard}, logically and symbolically oriented AI.
\item \textbf{Evolutionary computation} (EC) is a superset that also includes
  methods which are not evolutionary in the narrow sense (swarm algorithms,
  memetic algorithms, tabu search).
\item \textbf{Evolutionary algorithms} (EA) are the subset of EC that uses the
  basic principles of natural evolution: a population of candidate solutions,
  stochastic operators (crossover, mutation) and selection driven by an
  objective function.
\end{itemize}

EAs are \textbf{population-based stochastic search algorithms}. They are used
where no gradient or analytical form of the objective function is available
(so-called \emph{black-box optimisation}), where the solution space is discrete
or structured (trees, graphs, permutations), or where several mutually
non-dominated solutions are sought at once.

\subsubsection{Biological inspiration}

\begin{itemize}
\item \textbf{Darwin (1859, On the Origin of Species).} The environment has
  limited resources, reproduction is the key to survival, and better adapted
  individuals have a greater chance of reproducing. Successful phenotypic traits
  are reproduced, modified and recombined.
\item \textbf{Mendel (1866).} The gene as the basic unit of heredity; alleles are
  passed on independently. Reality is more complicated: \emph{polygeny} (several
  genes influence one trait), \emph{pleiotropy} (one gene influences several
  traits), mitochondrial DNA, epigenetics.
\item \textbf{Watson and Crick (1953).} The structure of DNA; the
  molecular-biological view of how information is stored and inherited. A codon
  (a triple of nucleotides) encodes one of the amino acids, and the code is
  redundant.
\item \textbf{Genotype $\to$ phenotype.} The mapping is one-way and complex
  (transcription DNA${\to}$RNA, translation RNA${\to}$protein).
  \emph{Lamarckism}, by contrast, assumes the existence of an inverse mapping,
  i.e.\ the inheritance of acquired characteristics. It is rejected in biology,
  but it is used as a technique in memetic algorithms
  (chapter~\ref{sec:lokalni}).
\end{itemize}

\begin{defbox}[Glossary: nature vs.\ algorithm]
\begin{tabular}{@{}ll@{}}
environment & optimisation problem \\
individual / chromosome & candidate solution, point of the search space \\
gene, allele & component of a solution and its value \\
genotype & internal representation (encoding) of an individual \\
phenotype & the \uv{outside} of an individual, on which quality is evaluated \\
fitness & measure of solution quality (objective function) \\
population & multiset of candidate solutions \\
generation & one iteration of the algorithm \\
\end{tabular}
\end{defbox}

\subsubsection{The general scheme of an evolutionary algorithm}

\begin{defbox}[Evolutionary algorithm]
An evolutionary algorithm is an iterative stochastic process that maintains a
population $P(t)$ of candidate solutions and repeatedly creates from it the
population $P(t+1)$ by means of \emph{parent selection}, \emph{recombination},
\emph{mutation} and \emph{environmental selection}, until a termination
condition is met.
\end{defbox}

\begin{algorithm}[!ht]
\caption{General scheme of an evolutionary algorithm}
\begin{algorithmic}[1]
\Require representation, fitness $f$, operators, parameters
\State $t \gets 0$; initialise the population $P(0)$ at random
\State evaluate all individuals in $P(0)$ with the function $f$
\While{the termination condition is not met}
  \State \textbf{parent selection:} select a set of parents from $P(t)$ (probabilistically, according to fitness)
  \State \textbf{recombination:} from pairs ($n$-tuples) of parents create offspring with probability $p_C$
  \State \textbf{mutation:} perturb each offspring at random with probability $p_M$
  \State evaluate the offspring $P'(t+1)$
  \State \textbf{environmental selection:} from $P(t) \cup P'(t+1)$ (or from $P'(t+1)$ only) select the new population $P(t+1)$
  \State $t \gets t+1$
\EndWhile
\Ensure the best individual found
\end{algorithmic}
\end{algorithm}

This cycle is common to all EAs; the individual variants (GA, ES, EP, GP, DE)
differ in the choice of representation, in the emphasis placed on the individual
operators and in the type of selection. Two opposing forces play the key role:

\begin{itemize}
\item \textbf{Variation} (crossover and mutation) creates new diversity --- it
  provides \emph{exploration} of the space.
\item \textbf{Selection} removes diversity and steers the search towards
  promising regions --- it provides \emph{exploitation}.
\end{itemize}

The balance between them is the central problem of EA design. Exploration alone
degenerates into a random walk; exploitation alone leads to getting stuck in a
local optimum and to \emph{premature convergence}.

\subsubsection{Components of an evolutionary algorithm}

\paragraph{Representation.} The choice of encoding is the most important design
decision. The classical options are: a binary string, a vector of integers, a
vector of real numbers, a permutation, a tree, a graph, a matrix, a finite
automaton, a neural network. The representation determines which operators make
sense.

\paragraph{Fitness function.} A mapping $f:$ phenotype space $\to \mathbb{R}$
that measures quality. Minimisation problems are usually converted to
maximisation (e.g.\ $f' = C_{\max} - f$ or $f' = 1/(1+f)$). Fitness may be
\emph{deterministic}, \emph{noisy} (simulation), \emph{dynamic} (changing in
time, section~\ref{sec:koevoluce}) or \emph{contextual} (dependent on the other
individuals --- coevolution).

\paragraph{Population.} A multiset of individuals of (usually) fixed size $n$.
The population is the carrier of \emph{diversity}; diversity is measured, for
example, by the average Hamming distance, by the entropy of the alleles at the
individual positions, or by the variance of fitness.

\paragraph{Initialisation.} Usually uniformly random. The alternatives are
\emph{seeding} (inserting heuristic or expert solutions --- with the risk of
premature convergence), constructive heuristics (e.g.\ nearest neighbour for the
TSP), and Latin hypercubes for real-valued vectors.

\paragraph{Termination condition.} A maximum number of generations or fitness
evaluations, reaching a target quality, stagnation of the best fitness for $k$
generations, or loss of population diversity.

\subsubsection{Selection mechanisms}\label{ssec:selekce}

\paragraph{Roulette wheel (fitness-proportionate) selection}

\begin{defbox}[Roulette wheel selection]
The individual $x_i$ is selected with probability
\[
p_i \;=\; \frac{f(x_i)}{\sum_{j=1}^{n} f(x_j)} \;=\; \frac{f(x_i)}{n\,\bar{f}},
\]
where $\bar f$ is the average fitness of the population. The expected number of
times an individual is selected in $n$ draws is $n p_i = f(x_i)/\bar f$. This is
sampling \emph{with replacement} --- one individual may be selected several
times.
\end{defbox}

\textbf{Example (Goldberg's classical one).}\label{ex:goldberg} A population of
four 5-bit strings, fitness $f(x) = x^2$ (decoded as an integer):

\begin{center}
\begin{tabular}{@{}llrrrr@{}}
\toprule
$i$ & string & $x$ & $f=x^2$ & $p_i = f/\sum f$ & exp.\ count $f/\bar f$ \\
\midrule
1 & 01101 & 13 & 169 & 0.144 & 0.58 \\
2 & 11000 & 24 & 576 & 0.492 & 1.97 \\
3 & 01000 &  8 &  64 & 0.055 & 0.22 \\
4 & 10011 & 19 & 361 & 0.309 & 1.23 \\
\midrule
  & \multicolumn{2}{l}{sum} & 1170 & 1.000 & 4.00 \\
  & \multicolumn{2}{l}{average $\bar f$} & 292.5 & & \\
\bottomrule
\end{tabular}
\end{center}

The roulette wheel divides the circumference into sectors of sizes $p_i$ and is
spun $n$ times. Individual 2 will occupy roughly two places in the mating pool,
individual 3 will probably drop out.

\textbf{Disadvantages of roulette wheel selection:}
\begin{itemize}
\item \emph{Premature dominance} of a \uv{superman}: an individual with extreme
  fitness takes over the population and diversity disappears.
\item \emph{Loss of selection pressure} towards the end of the run: when all
  fitness values are similar, selection approaches a random draw.
\item It requires positive fitness values and is sensitive to affine
  transformations ($f$ and $f + c$ give different probabilities). The remedy is
  \emph{scaling} (section~\ref{ssec:skalovani}).
\item High variance of sampling --- the actual counts may differ considerably
  from the expected ones.
\end{itemize}

\paragraph{Stochastic universal sampling (SUS)}

Stochastic universal sampling uses a single spin of the wheel with $n$ evenly
spaced pointers at a distance of $1/n$. The expected numbers of selections
remain the same as for the roulette wheel, but the variance is minimal: an
individual with expected count $e_i$ is selected $\lfloor e_i \rfloor$ or
$\lceil e_i \rceil$ times. It is therefore a \uv{fairer roulette wheel}.

\paragraph{Tournament selection}

\begin{defbox}[$k$-tournament selection]
Select $k$ individuals at random (uniformly, with replacement) and return the
best of them. The parameter $k$ (typically 2, 3, 5) directly controls the
selection pressure: the probability that the $i$-th best individual out of $n$
is selected (\emph{here} $i=1$ denotes the \textbf{best} one) is
$\big((n-i+1)^k - (n-i)^k\big)/n^k$ --- it is the probability that all $k$ draws
fall among the $i$-th individual and the worse ones, minus the probability that
all $k$ fall strictly below the $i$-th. A deterministic tournament can be
softened: the better individual wins with probability $p < 1$.
\end{defbox}

Advantages: it requires no global information (no summation of fitness), it is
invariant with respect to monotone transformations of fitness, it is easy to
parallelise, and it works even where fitness is not given explicitly and the
comparison is obtained from \emph{playing} (a played match, a simulation,
coevolution).

\paragraph{Rank-based selection}

Rank-based selection replaces fitness by the rank of the individual. Linear rank
selection assigns to the individual of rank $i$ (\emph{note, here the ordering is
the opposite of the tournament case:} worst $i=1$, best $i=n$) the probability
\[
p_i \;=\; \frac{1}{n}\Big(s^{-} + (s^{+}-s^{-})\,\frac{i-1}{n-1}\Big),
\qquad 1 \le s^{+} \le 2,\quad s^{-} = 2 - s^{+},
\]
where $s^{+}$ is the expected number of copies of the best individual. The
exponential variant uses $p_i \propto (1-c)c^{\,n-i}$. Rank-based selection
maintains a constant selection pressure throughout the whole run and is robust
both against \uv{supermen} and against a change of the fitness scale.

\paragraph{Boltzmann selection}

The selection probability is $p_i \propto e^{f(x_i)/T}$ with a
\emph{temperature} $T$ that decreases in time: at the beginning (high $T$)
selection is almost uniform (exploration), at the end (low $T$) almost
deterministic (exploitation). Analogously the \emph{Boltzmann tournament}: the
worse individual wins with a probability depending exponentially on the ratio of
the fitness difference to the temperature. There is a direct link to simulated
annealing (chapter~\ref{sec:lokalni}).

\paragraph{Environmental selection, elitism, population models}

\begin{itemize}
\item \textbf{Generational model}: the whole population is replaced by the
  offspring, $|P'| = |P| = n$.
\item \textbf{Steady-state model}: in each step only $\lambda$ (often 1--2)
  offspring are created, and they replace selected individuals (the worst one, a
  random one, the most similar one --- \emph{crowding}). The parameter
  \emph{generation gap} $G = \lambda/n$ gives the proportion of the population
  that is replaced.
\item \textbf{Elitism}: the $k$ best individuals are copied unconditionally into
  the next generation. This guarantees monotonicity of the best fitness found
  and is a necessary condition for a number of convergence theorems; too strong
  an elitism, however, reduces diversity.
\item \textbf{Deterministic ES selection}: $(\mu+\lambda)$ selects the $\mu$ best
  individuals from parents and offspring together (implicit elitism),
  $(\mu,\lambda)$ from the offspring only (chapter~\ref{sec:es}).
\end{itemize}

\begin{defbox}[Selection pressure and takeover time]
\emph{Selection pressure} is the degree to which better individuals are
favoured. It is quantified by the \emph{takeover time} $\tau^*$ --- the number of
generations in which a single copy of the best individual (without variation)
fills the whole population. For a tournament of size $k$ we roughly have
$\tau^* \approx \ln n/\ln k$; for fitness-proportionate selection the takeover
time is longer and depends on the distribution of fitness.
\end{defbox}

\subsubsection{Theoretical limits: No Free Lunch}

\begin{theorem}[No Free Lunch, Wolpert \& Macready 1995--1997]
Averaged over \emph{all} possible objective functions on a finite domain, any
two deterministic non-repeating search algorithms have the same expected
performance. There is therefore no universally best optimisation algorithm.
\end{theorem}

The practical consequence: it pays to build \textbf{domain-specific variants} of EAs ---
a suitable representation and operators that respect the structure of the problem. The
theorem does not say that all algorithms are equally good on \emph{real} problems; real
functions form a very small and highly structured subset of all functions.

\subsubsection{Historical branches of evolutionary algorithms}

\begin{center}
\small
\begin{tabular}{@{}>{\raggedright\arraybackslash}p{2.3cm}
                  >{\raggedright\arraybackslash}p{3.1cm}
                  >{\raggedright\arraybackslash}p{3.1cm}
                  >{\raggedright\arraybackslash}p{3.1cm}
                  >{\raggedright\arraybackslash}p{3.1cm}@{}}
\toprule
 & \textbf{GA} & \textbf{ES} & \textbf{EP} & \textbf{GP} \\
\midrule
author, year & Holland, 1975 & Rechenberg, Schwefel, 1964 &
  L.~Fogel, Owens, Walsh, 1965 & Koza, 1989--92 \\[2pt]
representation & binary string & vector $\mathbb{R}^n$ + strategy parameters &
  finite automaton (genotype = phenotype) & syntax tree (LISP) \\[2pt]
main operator & crossover & mutation & mutation & subtree crossover \\[2pt]
crossover & 1-point & multi-parent recombination (local/global) &
  \emph{none} & subtree exchange \\[2pt]
mutation & bit-flip $p_M$ & Gaussian, self-adaptive $\sigma$ &
  \uv{clever} automaton mutations & several types \\[2pt]
parent selection & roulette wheel & random & each one once & tournament \\[2pt]
env.\ selection & generational replacement &
  deterministic $(\mu,\lambda)$, $(\mu+\lambda)$ &
  tournament (parents $+$ offspring) & generational \\[2pt]
theory & schemata & convergence, 1/5 rule & --- & bloat, schemata \\
\bottomrule
\end{tabular}
\end{center}

Today the boundaries between the branches are blurred (EP and ES, for instance,
have practically merged) and we speak uniformly of evolutionary algorithms.

\subsection{Genetic algorithms}
\label{ssec:sga}
\subsubsection{The simple genetic algorithm (SGA)}

Holland's original 1975 proposal is nowadays referred to as the \emph{simple
genetic algorithm} (SGA). It uses a minimal set of operators and the simplest
encoding, and it was designed so as to make theoretical analysis possible
(schema theory, chapter~\ref{sec:schemata}).

\begin{defbox}[SGA]
A population of $n$ binary strings of length $m$; fitness-proportionate
(roulette wheel) selection; one-point crossover with probability $p_C$; bit
mutation with probability $p_M$ per bit; generational replacement of the
population ($P(t+1)$ replaces $P(t)$ entirely).
\end{defbox}

\begin{algorithm}[!ht]
\caption{Simple genetic algorithm (SGA)}
\begin{algorithmic}[1]
\State $t \gets 0$; generate $P(0)$ = $n$ strings of length $m$ at random
\While{the termination condition is not met}
  \State compute $f(x)$ for every $x \in P(t)$
  \State $P(t+1) \gets \emptyset$
  \For{$n/2$ times}
    \State select a pair of parents $x, y$ from $P(t)$ by the roulette wheel
    \State with probability $p_C$ cross $x, y$ by one-point crossover
    \State flip every bit of $x$ and $y$ with probability $p_M$
    \State insert $x, y$ into $P(t+1)$
  \EndFor
  \State $t \gets t+1$
\EndWhile
\end{algorithmic}
\end{algorithm}

Typical parameter values: $n \in [50, 500]$, $p_C \in [0.6, 0.9]$,
$p_M \approx 1/m$ (i.e.\ on average one bit per individual is changed).

\subsubsection{Binary representation}

\paragraph{Encoding of integers and real numbers.} A real number $x \in [a,b]$
encoded on $k$ bits as an integer $z \in \{0,\dots,2^k-1\}$ is decoded as
\[
x \;=\; a \;+\; z \cdot \frac{b-a}{2^k - 1}.
\]
The precision is $(b-a)/(2^k-1)$. Advantage: explicit control over the precision
and compression of the search space. Disadvantage: \emph{Hamming cliffs} ---
neighbouring values may differ in all bits (0111 $\to$ 1000), so a small change
of the phenotype requires a large change of the genotype.

\paragraph{Gray code.} It solves the Hamming cliffs: neighbouring integers
differ in exactly one bit in the Gray code. Conversion from binary $b$ to Gray
$g$: $g_1 = b_1$, $g_i = b_{i-1} \oplus b_i$. Back: $b_1 = g_1$,
$b_i = b_{i-1} \oplus g_i$.

\paragraph{Holland's argument for the binary alphabet.} Binary strings of length
100 encode roughly the same information as decimal strings of length 30, but
they contain more schemata ($3^{100}$ against $11^{30}$), so a GA
\uv{processes} more information. Today we know that schemata are not as
fundamental as Holland believed, and that what matters is the
\textbf{naturalness of the encoding for the given problem}. Real-valued vectors
are used for real-valued optimisation, permutations for the TSP, and so on.

\subsubsection{Crossover operators}

\begin{defbox}[Crossover (recombination)]
An operator that creates offspring from two (or more) parents by combining their
genetic material. In a GA it is the main operator; it expresses the hope that
recombining good parts (\emph{building blocks}) will produce even better
individuals. It is applied with probability $p_C$, otherwise the offspring are
copies of the parents.
\end{defbox}

\paragraph{One-point crossover.} Choose a cut point $k \in \{1,\dots,m-1\}$ at
random and exchange the trailing parts.

\textbf{Example.} Parents $P_1 = 1011|0110$, $P_2 = 0010|1101$, cut point $k=4$:
\[
O_1 = \mathbf{1011}\,1101, \qquad O_2 = \mathbf{0010}\,0110.
\]

\paragraph{Two-point crossover.} Choose two points $k_1 < k_2$ and exchange the
middle segment. The advantage over one-point crossover: the string behaves like
a ring, so genes at the ends are not disadvantaged (with one-point crossover the
probability of disrupting a segment is proportional to its \emph{defining
length}, see chapter~\ref{sec:schemata}).

\textbf{Example.} $P_1 = 10|1101|10$, $P_2 = 00|1011|01$ $\Rightarrow$
$O_1 = 10\,1011\,10$, $O_2 = 00\,1101\,01$.

\paragraph{Uniform crossover.} For each position $j$ a coin is tossed: with
probability $0.5$ the gene comes from the first parent, otherwise from the
second one (the second offspring is complementary). It does not favour adjacent
positions, but it strongly disrupts schemata with a large defining length ---
it has a high \emph{disruptiveness}.

\textbf{Example.} $P_1 = 11111111$, $P_2 = 00000000$, mask $10110010$:
$O_1 = 10110010$, $O_2 = 01001101$.

\begin{center}
\begin{tabular}{@{}lll@{}}
\toprule
\textbf{operator} & \textbf{prob.\ of disrupting a schema} & \textbf{note} \\
\midrule
one-point & $d/(m-1)$ & positional bias, protects short schemata \\
two-point & $\dfrac{2d(m-d)}{m(m-1)}$ & string as a ring, no end bias \\
$k$-point & grows with $k$ & \\
uniform & $1 - (1/2)^{o(S)-1}$ & distributional bias, most disruptive \\
\bottomrule
\end{tabular}
\end{center}

\textbf{A remark on two-point crossover.} If we regard the string as a ring, we
choose two of the $m$ cuts; a schema is disrupted if \emph{exactly one} cut
falls inside its defining segment of $d$ gaps, i.e.\ with probability
$d(m-d)/\binom{m}{2} = 2d(m-d)/(m(m-1))$. For a schema stretching from the first
to the last bit ($d = m-1$) the disruption drops from certainty (one-point:
$d/(m-1) = 1$) to a mere $2/m$ --- this is exactly the \uv{removal of the end
bias}. For short schemata, on the contrary, two-point crossover is about twice
as destructive as one-point crossover.

\subsubsection{Mutation and inversion}

\paragraph{Bit mutation.} Every bit is flipped independently with probability
$p_M$. In the SGA mutation is a \emph{secondary} operator --- it acts as an
insurance against getting stuck in a local optimum and against the permanent
loss of an allele from the population. (In EP and in the early ES, by contrast,
mutation is the only source of variability.) Recommendation:
$p_M \approx 1/m$.

\paragraph{Inversion.} Holland's original proposal also contained the operator
of \emph{inversion}: it reverses the order of a part of the string but
\emph{preserves the meaning of the bits} (a gene carries its own position
identifier). The goal was to make it possible to evolve the ordering of the
genes itself, so that cooperating genes lie close to one another and are not
disrupted by crossover. This is technically demanding and the benefit was never
demonstrated; with a permutation representation, however, inversion is used as a
very effective mutation (section~\ref{sec:reprez}).

\subsubsection{Fitness scaling}\label{ssec:skalovani}

Roulette wheel selection is sensitive to the absolute values of fitness. At the
beginning of a run the variance of fitness tends to be large (a few
\uv{supermen} take over the population), at the end small (selection loses
pressure). The remedy is a transformation $f \mapsto f'$:

\begin{itemize}
\item \textbf{Linear scaling:} $f' = a f + b$, where $a, b$ are chosen so that
  $\bar{f'} = \bar{f}$ and $f'_{\max} = C_{\mathrm{mult}} \cdot \bar{f}$,
  typically $C_{\mathrm{mult}} \in [1.2, 2]$. A safeguard against negative
  values for the worst individuals is necessary.
\item \textbf{Sigma truncation:} $f' = \max\big(0,\; f - (\bar f - c\,\sigma_f)\big)$,
  where $\sigma_f$ is the standard deviation of fitness and $c \in [1,3]$. It
  removes outlying low values and adapts automatically to the variance.
\item \textbf{Power scaling:} $f' = f^{\,k}$, where $k$ is increased during the
  run (the pressure grows).
\item \textbf{Rank scaling:} replaces fitness by rank (see
  section~\ref{ssec:selekce}) --- the most robust and nowadays the most common
  option.
\item \textbf{Boltzmann scaling:} $f' = e^{f/T}$ with a decreasing temperature.
\end{itemize}

\textbf{Example.} A population with fitness values $(169, 576, 64, 361)$,
$\bar f = 292.5$, $f_{\max} = 576$. Linear scaling with
$C_{\mathrm{mult}} = 1.5$: we want $f'_{\max} = 438.75$ and $\bar{f'} = 292.5$.
From the conditions $a f_{\max} + b = 438.75$ and $a \bar f + b = 292.5$ we get
$a = 146.25/283.5 = 0.516$, $b = 292.5 - 0.516 \cdot 292.5 = 141.6$.
The new fitness values are $(228.8;\, 438.8;\, 174.6;\, 327.9)$; the share of
the best individual dropped from $49.2\,\%$ to $37.5\,\%$ --- the selection
pressure is moderated and the worst individual does not lose its chance.

\subsection{Representation of an individual and genetic operators}
\label{sec:reprez}
The representation determines which operators make sense, and thereby also what
the \emph{fitness landscape} looks like --- that is, how easy the search is.
This is the most important practical decision when designing an EA.

\subsubsection{Integer representation}

An individual is a vector $(z_1,\dots,z_m) \in \prod_j D_j$, where the $D_j$ are
finite domains. We distinguish:
\begin{itemize}
\item \textbf{Cardinal (ordinal) attributes} --- a meaningful ordering and
  metric exist on the values (e.g.\ the number of machines, age).
\item \textbf{Nominal attributes} --- the values are mere codes (colour, type of
  component). Here it makes no sense to speak of \uv{close} values.
\end{itemize}

\paragraph{Mutation.}
\begin{itemize}
\item \emph{Unbiased} (\emph{random resetting}): the new value is drawn at
  random from the whole domain. The only correct choice for nominal attributes.
\item \emph{Biased} (\emph{creep mutation}): the new value is the original one
  shifted by a small random step (e.g.\ from a discretised normal
  distribution). Suitable only for ordinal attributes.
\end{itemize}

\paragraph{Crossover.} One-point, multi-point, uniform --- exactly as with the
binary representation, only integers are copied instead of bits.

\subsubsection{Real-valued representation}

An individual is $\vec{x} = (x_1,\dots,x_n) \in \mathbb{R}^n$. Historically real
numbers were encoded into bit strings; today one works directly with real values
and the operators are adapted accordingly.

\paragraph{Mutation.}
\begin{itemize}
\item \textbf{Gaussian (normal) mutation:}
  $x_i' = x_i + \sigma_i \cdot N(0,1)$. The parameter $\sigma$ (the step width)
  is crucial; it may be fixed, deterministically decreasing
  (e.g.\ $\sigma(t) = 1 - 0.9\,t/T$), adaptive (the 1/5 rule) or self-adaptive
  (part of the genome, see chapter~\ref{sec:es}).
\item \textbf{Cauchy mutation:} heavier tails $\Rightarrow$ more frequent large
  jumps $\Rightarrow$ better escape from local optima (used in \emph{fast EP}).
\item \textbf{Uniform mutation:} $x_i'$ drawn at random from $[a_i, b_i]$ ---
  unbiased.
\item Handling of the bounds: clipping, reflection, periodic wrapping or
  resampling.
\end{itemize}

\paragraph{Crossover.} Two types are distinguished:
\begin{itemize}
\item \textbf{Structural (discrete):} 1-point, uniform --- the values are merely
  rearranged between the parents; the offspring is a vertex of the hyperbox
  given by the parents.
\item \textbf{Arithmetic:} the values are combined.
\end{itemize}

\begin{defbox}[Arithmetic crossover]
For parents $\vec x, \vec y$ and $\alpha \in (0,1)$:
\[
\vec z \;=\; \alpha \vec x + (1-\alpha)\vec y \qquad\text{(convex combination)}.
\]
Variants: \emph{whole} arithmetic crossover (all components are crossed --- the
most common one), \emph{simple arithmetic} (only one randomly chosen
component), \emph{single} (a combination with one-point crossover: beyond the
cut point the values are combined arithmetically). For $\alpha = 1/2$ it is the
plain average of the parents.
\end{defbox}

\textbf{Example.} $\vec x = (2.0;\, -1.0;\, 4.0)$,
$\vec y = (0.0;\, 3.0;\, 1.0)$, $\alpha = 0.25$:
\[
\vec z = 0.25\,\vec x + 0.75\,\vec y = (0.5;\; 2.0;\; 1.75).
\]

The disadvantage of purely arithmetic crossover: the offspring lie inside the
convex hull of the parents, so the variance of the population necessarily
decreases (\uv{shrinking}). Extended variants are therefore used:
\begin{itemize}
\item \textbf{BLX-$\alpha$} (\emph{blend crossover}): $z_i$ is drawn uniformly
  from the interval $[\min - \alpha I,\; \max + \alpha I]$, where
  $I = |x_i - y_i|$; typically $\alpha = 0.5$. It allows the offspring to lie
  outside the parents as well.
\item \textbf{SBX} (\emph{simulated binary crossover}): it imitates the
  behaviour of one-point binary crossover in a real-valued space. The offspring
  arise by \uv{stretching} the pair of parents around their midpoint by a
  factor $\beta$:
  \[
  z_{1,i} = \tfrac12\big[(1+\beta)x_i + (1-\beta)y_i\big],\qquad
  z_{2,i} = \tfrac12\big[(1-\beta)x_i + (1+\beta)y_i\big],
  \]
  so the midpoint of the parents is preserved and $|z_1 - z_2| = \beta\,|x -
  y|$. The factor is drawn from $u \sim U(0,1)$: $\beta = (2u)^{1/(\eta+1)}$
  for $u \le 1/2$, otherwise $\beta = \big(2(1-u)\big)^{-1/(\eta+1)}$.
  Intuition: $\beta = 1$ returns the parents, $\beta < 1$ gives offspring
  between them, $\beta > 1$ outside; a large $\eta$ pushes $\beta$ towards one
  (offspring close to the parents), a small $\eta$ allows large jumps. The
  standard operator in NSGA-II.
\end{itemize}

\subsubsection{Permutation representation}

An individual is a permutation $\pi$ of $\{1,\dots,n\}$. Typical problems: the
travelling salesman problem (TSP), scheduling, assignment problems, ordering of
operations. The classical operators fail here --- the result of one-point
crossover of two permutations is in general not a permutation. Operators are
therefore needed that \emph{preserve feasibility} and at the same time inherit
something meaningful. Which information should be inherited depends on the
problem:

\begin{center}
\begin{tabular}{@{}lll@{}}
\toprule
\textbf{operator} & \textbf{what it preserves} & \textbf{typical use} \\
\midrule
PMX & absolute positions of the cities & problems where position matters \\
OX & relative order & TSP (cyclic tours) \\
CX & absolute positions (cycles) & problems where position matters \\
ER & edges (adjacencies) & TSP --- the best quality \\
\bottomrule
\end{tabular}
\end{center}

\paragraph{PMX --- partially mapped crossover}

\begin{defbox}[PMX (partially mapped crossover, Goldberg)]
A two-point operator. (1)~Choose the segment between the two cut points and
exchange it between the parents. (2)~From the exchanged segments derive the
\emph{mapping} of the corresponding values. (3)~Copy the remaining positions
from the original parent; where a conflict would arise (the value is already in
the offspring), use the mapping (repeatedly if necessary).
\end{defbox}

\textbf{Example.} $P_1 = (123|4567|89)$, $P_2 = (452|1876|93)$.
\begin{enumerate}
\item Exchange of the segments: $O_1 = (\ldots|1876|\ldots)$,
  $O_2 = (\ldots|4567|\ldots)$.
\item Mapping from the segments: $1{\leftrightarrow}4$, $8{\leftrightarrow}5$,
  $7{\leftrightarrow}6$, $6{\leftrightarrow}7$.
\item Filling in the conflict-free positions from $P_1$ into $O_1$:
  $O_1 = (.23|1876|.9)$ (the values 1 and 8 would collide).
\item The conflicts are resolved by the mapping: $1 \to 4$, $8 \to 5$, hence
  $O_1 = (423|1876|59)$. Analogously $O_2 = (182|4567|93)$.
\end{enumerate}

\paragraph{OX --- order crossover}

\begin{defbox}[OX (order crossover, Davis)]
A two-point operator preserving the \emph{relative order}. (1)~Copy the segment
between the cut points from the first parent. (2)~From the second parent read
the values cyclically, starting at the position after the second cut point, and
skip those that are already in the offspring. (3)~Insert the remaining values
into the free positions of the offspring, again cyclically from the position
after the second cut point.
\end{defbox}

\textbf{Example.} $P_1 = (123|4567|89)$, $P_2 = (452|1876|93)$.
\begin{enumerate}
\item $O_1$ takes over the segment from $P_1$:
  $O_1 = ({\cdot}{\cdot}{\cdot}|4567|{\cdot}{\cdot})$.
\item From $P_2$ we read cyclically from position 8: $9,3,4,5,2,1,8,7,6$.
\item We delete the values already contained ($4,5,6,7$): $9,3,2,1,8$ remain.
\item We fill in from position 8 cyclically (positions $8,9,1,2,3$):
  $O_1 = (218|4567|93)$.
\item Symmetrically (segment from $P_2$, order from $P_1$):
  $O_2 = (345|1876|92)$.
\end{enumerate}
The principle: the inner segment is taken over unchanged, the rest is filled in
in the relative order of the second parent, skipping duplicates. For the TSP, OX
is more suitable than PMX, because in a round tour it is not the absolute
position of a city that matters, but the order.

\paragraph{CX --- cyclic crossover}

\begin{defbox}[CX (cyclic crossover, Oliver)]
It preserves the \emph{absolute position}: every gene of the offspring comes
from the same position in one of the parents. Cycles are constructed: start at
position 1 and take the value from $P_1$; look at which value stands at the same
position in $P_2$, find it in $P_1$ and continue until the cycle closes. Take
the even cycles from one parent and the odd ones from the other.
\end{defbox}

\textbf{Example.} $P_1 = (123456789)$, $P_2 = (412876935)$.
We start with the value 1 from $P_1$ at position 1. Below it, in $P_2$, is the
value 4, which we find in $P_1$ at position 4 $\Rightarrow$ we add it; below it
is 8, which is in $P_1$ at position 8 $\Rightarrow$ we add it; below it is 3
(position 3), below that 2 (position 2), and below 2 is 1 --- the cycle has
closed. We have $O_1 = (1\,2\,3\,4\,\_\,\_\,\_\,8\,\_)$; the rest is filled in
from $P_2$: $O_1 = (123476985)$, and analogously $O_2 = (412856739)$.

\paragraph{ER --- edge recombination}

An observation: PMX, OX and CX preserve only about $60\,\%$ of the parents'
edges. For the TSP, however, it is precisely the \emph{edge} that carries the
quality. \emph{Edge recombination} (Whitley) therefore:
\begin{enumerate}
\item For every city it builds a \emph{neighbour list} from \textbf{both}
  parents.
\item It starts at a random (or the first) city.
\item The next city is the not-yet-visited neighbour that has the \emph{fewest}
  remaining neighbours (the heuristic being: do not leave cities
  \uv{stranded}); ties are broken at random.
\item If the neighbour list turns out to be empty, one continues with a random
  unvisited city (this happens in $1$--$1.5\,\%$ of the cases).
\end{enumerate}
\textbf{ER2} additionally marks the edges that occur in \emph{both} parents
(with the symbol \#) and prefers them when choosing.

\textbf{Example.} $P_1 = (123456789)$, $P_2 = (412876935)$ (cyclic tours).
Neighbour lists: $1{:}\,9,2,4$; $2{:}\,1,3,8$; $3{:}\,2,4,9,5$; $4{:}\,3,5,1$;
$5{:}\,4,6,3$; $6{:}\,5,7,9$; $7{:}\,6,8$; $8{:}\,7,9,2$; $9{:}\,8,1,6,3$.
Start at 1: the candidates are 9 (4~neighbours), 2 (3), 4 (3) --- 9 drops out,
between 2 and 4 we draw lots and 4 comes out. The neighbours of 4 are 3 and 5;
3 still has 3 free neighbours, 5 only 2, so we take 5, and so on --- the result
is e.g.\ $(145678239)$.

\emph{Technical note:} strictly speaking, the visited city is first removed from
\emph{all} the lists and only then are the remaining neighbours counted (above,
at the start, the counts are given \emph{before} removing the 1: after the
removal they would be 9$:$3, 2$:$2, 4$:$2). Here this changes nothing about the
order of the candidates, but when computing on paper one has to proceed this
way, otherwise different ties come out.

\paragraph{Mutation of permutations}

\begin{itemize}
\item \textbf{Swap:} exchange of two random cities.
\item \textbf{Insert:} removal of a city and its insertion elsewhere (the rest
  is shifted).
\item \textbf{Inversion (a 2-opt step):} reversal of a contiguous segment ---
  the most effective one for the TSP, because it changes only two edges.
\item \textbf{Scramble:} random shuffling of a contiguous segment.
\item \textbf{Exchange of subtours} and heuristic mutations (2-opt, 3-opt,
  Lin-Kernighan).
\end{itemize}

\subsubsection{Other representations}

Trees (section~\ref{sec:gp}), graphs and DAGs (Cartesian GP, neuroevolution),
matrices (schedules, the TSP as a predecessor matrix), finite automata (EP),
rule lists (LCS, section~\ref{sec:lcs}), artificial-life agents.

\subsection{Genetic and evolutionary programming}
\label{sec:gp}
\subsubsection{History}

The idea of evolving programs was already mentioned by Alan Turing in the 1950s.
Forsyth (1980, the BEAGLE system) applied it to pattern recognition, Nichael
Cramer (1985) was the first to describe tree-shaped individuals, and John Koza
(1989--1992) formulated tree-based genetic programming in its present form
(including a patent).

\begin{defbox}[Genetic programming (GP)]
An evolutionary algorithm whose individuals are \textbf{executable structures}
--- programs, expressions, circuits, models. Fitness is determined by
\emph{running} the individual on test data. Neither the size nor the shape of an
individual is fixed; both change in the course of evolution.
\end{defbox}

\begin{algorithm}[!ht]
\caption{The general scheme of GP}
\begin{algorithmic}[1]
\State generate an initial population of random programs
\While{a sufficiently good solution has not been found}
  \State evaluate the programs by \textbf{running} them on the training data
  \State selection according to fitness (typically a tournament)
  \State crossover of two programs, mutation of programs
  \State create the new population
\EndWhile
\end{algorithmic}
\end{algorithm}

\subsubsection{Tree-based GP}

\paragraph{Representation.} A program is a \emph{syntax tree}. We define two
sets:
\begin{itemize}
\item \textbf{The terminal set} $T$ --- the leaves: variables, constants,
  functions without arguments (e.g.\ reading a sensor).
\item \textbf{The non-terminal (function) set} $F$ --- the inner nodes with a
  given \emph{arity}: arithmetic operations, logical connectives, conditionals,
  loops.
\end{itemize}
Two properties are required: \emph{closure} --- every function must be able to
process any output of any other one (hence the \uv{protected division} $\%$,
which returns 1 for a zero divisor) --- and \emph{sufficiency} --- the solution
must be expressible from $T \cup F$ at all.

Example: the tree \texttt{(+ (* x x) (sin y))} represents the expression
$x^2 + \sin y$.

\paragraph{Initialisation.}
\begin{itemize}
\item \textbf{Grow:} the nodes are drawn from $T \cup F$ until the limit on
  depth / number of nodes is reached $\Rightarrow$ irregular trees of various
  sizes.
\item \textbf{Full:} up to the given depth only $F$ is drawn from, at the last
  level only $T$ $\Rightarrow$ full, regular trees.
\item \textbf{Ramped half-and-half} (Koza): half of the population by the grow
  method, half by the full method, for uniformly distributed depths (e.g.\
  2--6). The standard choice; it produces rather \uv{bushy} regular shapes,
  which suits problems with symmetries where all the inputs are similarly
  important.
\item \textbf{Uniform sampling} of all possible trees gives more irregular
  shapes (some terminals closer to the root).
\item \textbf{Seeding:} inserting partial or expert solutions into the initial
  population --- it speeds things up, but premature convergence is a risk.
\end{itemize}

\paragraph{Crossover.}
\begin{itemize}
\item \textbf{Subtree exchange} (Koza's standard crossover): a random node is
  chosen in each parent and the subtrees below them are exchanged.
  Non-terminals tend to have a higher probability of being chosen (typically
  90~\% non-terminal, 10~\% terminal) --- otherwise only the leaves would often
  be exchanged.
\item \textbf{Homologous crossover} tries to preserve the position of the
  genetic material. First the \emph{common region} $C$ is found (by traversing
  from the root as long as the arities of the nodes agree):
  \begin{itemize}
  \item \emph{One-point:} the crossover point is chosen from the common region.
  \item \emph{Uniform (GPUX, Poli and Langdon):} every node of the common region
    is exchanged with a constant probability; nodes \emph{inside} $C$ are
    exchanged without any effect on their subtrees, nodes on the \emph{boundary}
    of $C$ are exchanged together with their subtrees. In the early phases large
    subtrees are thus exchanged, and as the population converges the operator
    becomes ever more local.
  \end{itemize}
\item \textbf{Context preserving crossover}: a node is identified by its
  \emph{coordinates} --- the sequence of branchings on the path from the root
  (e.g.\ $(2,1,3,1)$). The \emph{strong} variant allows crossover only between
  nodes with the \emph{same} coordinates, the \emph{weak} one is more
  permissive.
\end{itemize}

\paragraph{Mutation.} In GP it is necessary (not merely useful) to use
\emph{several types} of mutation:
\begin{itemize}
\item \textbf{Subtree mutation} --- replaces a random subtree by a new random
  one.
\item \textbf{Size-fair} subtree mutation --- the size of the new subtree is
  drawn so as to match the removed one.
\item \textbf{Node mutation} (node replacement) --- a node is exchanged for
  another one of the \emph{same arity}.
\item \textbf{Hoist} --- the individual is replaced by one of its subtrees
  (it shrinks).
\item \textbf{Shrink} --- a subtree is replaced by a terminal (it shrinks).
\item \textbf{Permutation mutation} --- the subtrees of one non-terminal are
  swapped.
\item \textbf{Mutation of constants} --- GP is traditionally bad at fine-tuning
  numerical values. Arithmetic mutation of constants helps, or even local
  optimisation of \emph{all} the constants in the tree (hill climbing, a
  gradient method) --- a memetic approach.
\end{itemize}

\paragraph{Fitness and selection.} Fitness = the error on the training data
(symbolic regression), the number of correct answers, the quality of the
behaviour in a simulation. Selection is usually by tournament. The risk of
overfitting is the same as in machine learning (it is handled by a validation
set and by penalising complexity).

\subsubsection{Bloat}

\begin{defbox}[Bloat]
\emph{Bloat} is the tendency of programs in GP to grow in size without a
corresponding improvement in fitness. Typically it is the accumulation of
\emph{introns} --- pieces of code that have no effect on the output
(e.g.\ \texttt{(+ x 0)}, \texttt{(if false ...)}).
\end{defbox}

Why it happens: in the space of programs there are far more large programs with the same
behaviour than small ones, so neutral drift alone pushes towards growth; moreover, introns
protect an individual against destructive crossover, so they \uv{hitchhike} along with the
good code. In nature growth runs into physical and energetic limits; in GP bloat has to be
fought explicitly:
\begin{itemize}
\item \textbf{Hard limits} on the depth or the number of nodes (an offspring exceeding the
  limit is discarded or replaced by the parent).
\item \textbf{A penalty in the fitness} (\emph{parsimony pressure}):
  $f' = f - \alpha \cdot \text{size}$; or a \emph{lexicographic} variant: with equal fitness
  the smaller individual wins.
\item \textbf{Anti-bloat operators} --- watch crossover and mutation so that they do not
  enlarge the individual much (size-fair variants), plus the shrinking mutations
  \emph{hoist} and \emph{shrink}.
\item \textbf{A multi-objective approach} --- size as a second criterion
  (section~\ref{sec:moea}).
\end{itemize}

\subsubsection{ADFs and typed GP}

\begin{defbox}[ADF (automatically defined functions)]
Subprograms evolved together with the main program --- a necessary condition for
\textbf{modularity}, higher complexity and the ability to capture the symmetries of the
problem. An individual consists of a \texttt{main} branch and one or more ADF branches;
every ADF is characterised by its \emph{arity} and may have a restricted (or different) set
of terminals and non-terminals. A call of an ADF behaves in the main program as a
\textbf{new non-terminal} of the given arity. The genetic operators work on the branches
\textbf{separately} (main is crossed with main, ADF$_1$ with ADF$_1$).

Extensions: automatically defined loops, iterations and recursion. An alternative is the
\emph{coevolution} of two populations --- of main programs and of subprograms
(section~\ref{sec:koevoluce}).
\end{defbox}

\textbf{Enforcing structure.} \emph{Strongly typed GP} gives every node a type of its inputs
and output and the operators may create only type-correct trees; \emph{grammar-based GP}
describes the admissible shapes of the trees by a context-free grammar. Both restrict the
search space to meaningful programs.

\subsubsection{Linear and graph GP}

\paragraph{Linear GP.} A program is a sequence of instructions (often machine
code or byte code) operating on registers. Advantages: simple operators (vector
crossover and mutation), fast interpretation, natural closeness to real
hardware. Disadvantage: a high risk of producing nonsensical programs. It is
popular in artificial life and in the evolution of game bots. A famous example:
\textbf{Tierra} (T.~Ray) --- a simulation of an artificial ecosystem in which
the individuals are self-copying machine-code programs competing for memory and
processor time; parasites, hyperparasites and defences against them emerged.
Successors: Avida, Avida-ED.

\paragraph{Graph GP.} A program is a general (often acyclic) graph. Originally
an extension of tree GP to parallel programs; it turned out that graphs describe
electrical circuits, finite automata, neural networks and plans well. The
problem is the complexity of the operators --- how to cross general graphs.

\paragraph{Cartesian GP (CGP).} \mimo{} An elegant way around this: a DAG is represented
by its planar image in a 2D grid with the coordinates of the nodes (the
phenotype), whereas the genotype is a \textbf{linear} sequence of integers
describing, for each node, its function (an index into a table of functions),
its inputs, and finally the output nodes of the graph. Properties:
\begin{itemize}
\item Mutation is a point mutation of the linear genome; one has to make sure
  that only valid functions and valid connections arise (a limit on the depth of
  the links).
\item A small change of the genotype may have a large impact on the phenotype;
  conversely, a number of mutations are \emph{neutral} (nodes not connected to
  the output form natural introns, which is beneficial for CGP).
\item Crossover is of little importance; a naive one-point crossover is too
  destructive, so an arithmetic variant and cuts at node boundaries are used.
\item Selection: usually a $(1+\lambda)$-ES with $\lambda = 4$.
\item The genome has a constant length (bloat is limited by the grid).
\end{itemize}

\paragraph{Developmental GP.} Instead of evolving a graph directly, one evolves
a \emph{tree-program that builds the graph step by step}. One starts with an
\uv{embryo} and applies operations for the structure (serial/parallel
connection, three-way splitting) and for the elements (insert a resistor, a
capacitor). Applications: Koza --- the evolution of electrical circuits (the
rediscovery of patented designs), Gruau --- \emph{cellular encoding} for neural
networks, AutoML --- the evolution of scikit-learn models.

\subsubsection{Grammatical evolution}

\begin{defbox}[Grammatical evolution (GE)]
The genotype is a \textbf{linear list of integers} (of variable length), the phenotype is a
program derived from a context-free grammar. A number from the genome determines which rule
is used to expand the leftmost non-terminal:
\[
\text{rule} \;=\; (\text{gene} \bmod \text{the number of rules for this non-terminal}).
\]
\end{defbox}

The use of the modulo makes the encoding \textbf{robust} (every genome is interpretable) and,
because the genome is linear, \textbf{the standard GA operators} can be used --- that is the
main practical advantage. The grammar moreover guarantees syntactic correctness. The drawback
is non-locality: a small change of the genotype (especially at the beginning) may change the
phenotype completely (the \emph{ripple effect}). If the derivation does not terminate, the
genome is read again from the beginning (\emph{wrapping}) with a limit on the number of wraps.

\subsubsection{Evolutionary programming}

\begin{defbox}[Evolutionary programming (EP)]
A branch of EAs (L.~Fogel, 1965 --- older than GAs) whose goal was to evolve
\uv{artificial intelligence}: an agent that predicts the state of the
environment and chooses an appropriate action. The agent is represented by a
\textbf{finite automaton}; \textbf{genotype = phenotype} (no special encoding).
\textbf{Crossover is not used}, all variability comes from mutations.
\end{defbox}

\paragraph{EP over automata.} The environment is a sequence of symbols of a
finite alphabet. The automaton receives the symbols on its input, its output is
compared with the next symbol of the sequence; fitness = the success of the
prediction (absolute error, mean squared error), often with a penalty for the
size of the automaton. Mutations: a change of an output symbol, a change of the
target state of a transition, the addition of a state, the removal of a state, a
change of the initial state. Half of the population is replaced by new automata.

\paragraph{A famous example --- predicting primes.} The goal is to predict the
sequence
\[
111010100010100010100010000010\ldots,
\]
where 1 means \uv{this number is a prime} (i.e.\ to learn the sieve of
Eratosthenes). One proceeds iteratively: first a short sequence is learned, then
it is extended by one symbol. Fitness: a point for every correctly guessed
symbol minus $0.01 \cdot N$ for the number of states. The result is instructive:
the automata kept simplifying until a one-state automaton generating a constant
0 emerged --- primes are sparse, after all, so the constant answer \uv{not a
prime} has a high success rate. A classic illustration of how evolution exploits
a badly designed fitness function.

\paragraph{EP today.}
\begin{itemize}
\item The encoding should be \emph{natural} for the problem (usually genotype =
  phenotype).
\item A set of \uv{clever}, domain-specific mutations; no crossover.
\item Parent selection: every individual is a parent exactly once.
\item Environmental selection: a tournament between parents and offspring.
\item EP over real-valued vectors contains --- just as ES does --- the mutation
  variances in the genome ($\sigma_j' = \sigma_j(1 + \alpha N(0,1))$, then
  $x_j' = x_j + \sigma_j' N(0,1)$). EP and ES have thus practically merged.
\end{itemize}

\subsubsection{Is crossover good for anything? The headless chicken experiment}

The traditional defence of crossover = the exchange of building blocks. The
criticism (typical of EP): crossover is just a strange large random mutation
which, moreover, does not respect the encoding.

\begin{defbox}[Headless chicken test (Jones, 1995)]
We compare standard crossover with \emph{random} crossover, in which an
individual is crossed with a \textbf{randomly generated} string. If the result
is the same or better, then the operator under study does not in fact work as
crossover but as a macromutation --- \uv{let us not call a headless chicken a
chicken, even though it has many chicken-like features}.
\end{defbox}

Interpretation: if clear building blocks exist, true crossover is better. If
random crossover wins, it means that in the given problem/encoding there are no
building blocks --- a bad encoding, a bad operator, or a problem unsuitable for
crossover.

\subsection{Exam summary}

\begin{itemize}
\item EA = population-based stochastic search; the cycle \emph{initialisation
  $\to$ evaluation $\to$ parent selection $\to$ recombination $\to$ mutation
  $\to$ evaluation $\to$ environmental selection}.
\item Variation (crossover, mutation) creates diversity, selection removes it
  and steers the search; designing an EA = finding the exploration/exploitation
  balance.
\item Selection: \textbf{roulette wheel} ($p_i = f_i/\sum f_j$, high variance,
  sensitive to scaling), \textbf{SUS} (same expected values, low variance),
  \textbf{tournament} (pressure controlled by $k$, invariant to fitness
  transformations, needs no explicit fitness), \textbf{rank-based} (constant
  pressure), \textbf{Boltzmann} (pressure grows with time).
\item Elitism guarantees monotonicity of the best fitness; the generational and
  steady-state models differ in the proportion of the population replaced.
\item No Free Lunch: there is no universally best algorithm $\Rightarrow$ the
  representation and the operators must be specialised to the domain.
\item The historical branches GA / ES / EP / GP and their characteristic choices
  (see the comparison table) are worth being able to list from memory.
\item SGA = binary strings + roulette wheel + 1-point crossover + bit-flip
  mutation + generational replacement. Parameters: $n$, $p_C \approx 0.7$,
  $p_M \approx 1/m$.
\item Binary encoding of real numbers: $x = a + z(b-a)/(2^k-1)$; the problem of
  Hamming cliffs is solved by the Gray code.
\item Crossover: 1-point (disrupts a schema with probability $d(S)/(m-1)$),
  2-point (no end bias), uniform (most disruptive, no positional bias).
\item Mutation is secondary in a GA --- it prevents the irreversible loss of
  alleles; inversion is a historical operator for evolving the ordering of
  genes.
\item Fitness scaling (linear, sigma truncation, power, rank, Boltzmann)
  repairs the weaknesses of fitness-proportionate selection: the dominance of
  supermen at the beginning and the loss of pressure at the end.
\item Representation $+$ operators $=$ fitness landscape. For nominal attributes
  only unbiased mutations, for ordinal ones biased (creep) mutations as well.
\item Real-valued representation: Gaussian mutation $x' = x + \sigma N(0,1)$;
  arithmetic crossover $\alpha \vec x + (1-\alpha)\vec y$ (it reduces the
  variance), hence BLX-$\alpha$ and SBX.
\item Permutations: PMX (positions + mapping), OX (relative order),
  CX (absolute positions by means of cycles), ER (edges, the best for the TSP).
  Be able to demonstrate PMX and OX on a concrete eight-element example!
\item Mutation of permutations: swap, insert, inversion (2-opt), scramble.
\item Tree GP (Koza): terminals + non-terminals, closure and sufficiency;
  initialisation by grow / full / ramped half-and-half; crossover = subtree
  exchange (non-terminals with a higher probability), several types of mutation
  including mutation of constants; fitness = running the program.
\item Bloat = growth in size without an improvement in fitness (introns); three
  explanations: replication accuracy, removal bias, crossover bias. Defences:
  size limits, penalties (parsimony pressure), anti-bloat operators (hoist,
  shrink, size-fair), a multi-objective approach.
\item ADFs = evolved subprograms (modularity); the operators work on the
  branches separately. Typed GP and grammar-based GP restrict the space to
  meaningful programs.
\item Variants: linear GP (byte code, Tierra), graph GP,
  Cartesian GP (a linear genotype $\to$ a 2D DAG, point mutation, a
  $(1+4)$-ES), developmental GP (a tree builds a graph --- circuits, networks).
\item Grammatical evolution: a linear genome of integers, the choice of a
  grammar rule by means of the modulo, a derivation tree $\to$ the phenotype; a
  robust encoding and the standard GA operators, but a non-local mapping.
\item EP: finite automata, genotype = phenotype, mutation only, tournament
  selection from parents and offspring together; the example of predicting
  primes (a degenerate fitness).
\item The headless chicken experiment: a test of whether crossover really
  combines building blocks or is merely a macromutation.
\end{itemize}

%=====================================================================
\section{Schema theory and probabilistic models of the SGA}
\label{sec:schemata}
%=====================================================================

The chapter answers the question of \emph{why} genetic algorithms work. First Holland's
classical schema theory (and its criticism, which matters just as much at the examination as
the theorem itself), then the exact Vose model that superseded it.

\subsection{Schema theory}
\label{ssec:teorie-schemat}

Schema theory is the first attempt to explain formally \emph{why} genetic
algorithms work. It comes from Holland (1975) and was popularised by Goldberg
(1989). Today it is regarded as historically important but insufficient;
nevertheless it is a standard part of the examination, and its criticism is just
as important as the theorem itself.

\subsubsection{A schema, its order and defining length}

\begin{defbox}[Schema]
Let an individual be a word of length $m$ over the alphabet $\{0,1\}$. A
\textbf{schema} $S$ is a word of length $m$ over the alphabet $\{0,1,*\}$, where
$*$ means \uv{don't care}. A schema represents the set of all individuals that
agree with it in all fixed positions; a schema with $r$ stars represents $2^r$
individuals.
\end{defbox}

Example: the schema $1*0*$ represents $\{1000, 1001, 1100, 1101\}$; the schema
$**\ldots*$ represents the whole space.

\begin{defbox}[Order, defining length, fitness of a schema]
For a schema $S$ of length $m$ we define:
\begin{itemize}
\item \textbf{Order} $o(S)$ = the number of fixed positions (the number of zeros
  and ones).
\item \textbf{Defining length} $d(S)$ = the distance between the first and the
  last fixed position ($d(S) = 0$ for a schema with a single fixed position).
\item \textbf{Fitness of the schema} $F(S,t)$ = the average fitness of those
  individuals in the population $P(t)$ that represent the schema $S$.
  \emph{Note:} this value depends on the current population, not only on $S$
  and $f$.
\end{itemize}
\end{defbox}

\textbf{Example.} For $S = 1{*}{*}0{*}1$ ($m=6$) we have $o(S) = 3$
(positions 1, 4, 6) and $d(S) = 6 - 1 = 5$.

\paragraph{Combinatorics of schemata.}
\begin{itemize}
\item There are $3^m$ schemata of length $m$ in total.
\item A single individual of length $m$ is a representative of $2^m$ schemata
  (at each position either its own bit or $*$).
\item A population of $n$ individuals represents between $2^m$ and
  $n \cdot 2^m$ schemata (depending on how similar they are).
\end{itemize}

\subsubsection{Holland's schema theorem}

\begin{theorem}[Schema theorem (TST), Holland 1975]
In a simple genetic algorithm with fitness-proportionate selection, one-point
crossover with probability $p_C$ and bit mutation with probability $p_M$, the
expected number of representatives of a schema $S$ in the next generation
satisfies
\[
\mathbb{E}\big[C(S,t+1)\big] \;\ge\; C(S,t)\,\frac{F(S,t)}{\bar f(t)}
\left[\,1 - p_C \frac{d(S)}{m-1} - p_M\, o(S)\right].
\]
In words: \textbf{short} (small defining length), \textbf{low-order} and
\textbf{above-average} schemata multiply exponentially over successive
generations.
\end{theorem}

\paragraph{Derivation}

Let $C(S,t)$ denote the number of individuals in $P(t)$ representing $S$, let
$n = |P(t)|$ and let $\bar f(t)$ be the average fitness of the population. We
proceed in three steps, one per operator.

\paragraph{Step 1 --- selection.} The probability of selecting an individual $v$
by the roulette wheel is $p_s(v) = f(v)/F(t)$, where
$F(t) = \sum_{u \in P(t)} f(u)$. The probability that the selected individual
represents $S$ is
\[
p_s(S) \;=\; \sum_{v \in S \cap P(t)} \frac{f(v)}{F(t)}
\;=\; \frac{C(S,t)\,F(S,t)}{F(t)}.
\]
For $n$ independent draws we therefore have
\[
\mathbb{E}\big[C(S,t+1)\big] \;=\; n\, p_s(S) \;=\; C(S,t)\,\frac{F(S,t)}{\bar f(t)},
\qquad \bar f(t) = F(t)/n .
\]
If the schema is above average by $\varepsilon$, i.e.\
$F(S,t) = \bar f(t)\,(1+\varepsilon)$, and this ratio is maintained in time, we
get
\[
C(S,t+1) = C(S,t)(1+\varepsilon)
\qquad\Longrightarrow\qquad
C(S,t) = C(S,0)\,(1+\varepsilon)^{t},
\]
that is, \textbf{exponential growth} (or decay for $\varepsilon<0$).

\paragraph{Step 2 --- crossover.} One-point crossover chooses the cut point from
$m-1$ possibilities. The schema is \emph{certainly} disrupted if the cut falls
between its first and its last fixed position, i.e.\ with probability at most
$d(S)/(m-1)$ (at most, because even with a cut inside, the schema may be
restored if the other parent contains the same values). Crossover is performed
with probability $p_C$, so
\[
p_{\mathrm{surv}}^{\mathrm{cross}}(S) \;\ge\; 1 - p_C\,\frac{d(S)}{m-1}.
\]

\paragraph{Step 3 --- mutation.} The schema survives if none of its $o(S)$ fixed
positions is changed:
\[
p_{\mathrm{surv}}^{\mathrm{mut}}(S) = (1-p_M)^{o(S)} \;\approx\; 1 - p_M\,o(S)
\quad \text{for } p_M \ll 1 .
\]

\paragraph{Putting it together.} Assuming (approximate) independence of the two
destructive events and neglecting their product:
\[
\mathbb{E}\big[C(S,t+1)\big] \;\ge\; C(S,t)\,\frac{F(S,t)}{\bar f(t)}
\left[1 - p_C \frac{d(S)}{m-1} - p_M o(S)\right]. \qquad\blacksquare
\]

\paragraph{Numerical example}

Take Goldberg's population from p.~\pageref{ex:goldberg} ($m=5$, $n=4$):
$01101\,(169)$, $11000\,(576)$, $01000\,(64)$, $10011\,(361)$;
$\bar f = 292.5$. The parameters are $p_C = 0.7$, $p_M = 0.001$.

\begin{center}
\begin{tabular}{@{}lccccc@{}}
\toprule
schema $S$ & representatives & $F(S)$ & $o(S)$, $d(S)$ & survival & $\mathbb{E}[C(S,t{+}1)]$ \\
\midrule
$1{*}{*}{*}{*}$ & 11000, 10011 & $468.5$ & $1;\ 0$ &
  $1 - 0 - 0.001 = 0.999$ &
  $2 \cdot 1.602 \cdot 0.999 = 3.2$ \\
$1{*}{*}{*}1$ & 10011 & $361$ & $2;\ 4$ &
  $1 - 0.7 - 0.002 = 0.298$ &
  $1 \cdot 1.234 \cdot 0.298 = 0.37$ \\
$0{*}{*}{*}{*}$ & 01101, 01000 & $116.5$ & $1;\ 0$ &
  $0.999$ &
  $2 \cdot 0.398 \cdot 0.999 = 0.80$ \\
\bottomrule
\end{tabular}
\end{center}

Interpretation: the short above-average schema $1{*}{*}{*}{*}$ spreads from 2 to
about 3.2 representatives. The schema $1{*}{*}{*}1$ is above average as well
($F/\bar f = 1.23$), but its large defining length is fatal --- crossover
disrupts it and we expect it to disappear. The below-average
$0{*}{*}{*}{*}$ declines. This is exactly what the theorem says.

\subsubsection{The building blocks hypothesis}

\begin{defbox}[Building blocks hypothesis (BBH)]
A genetic algorithm searches for a good solution by \textbf{recombining short,
low-order and above-average schemata} --- the so-called \emph{building blocks}
--- into ever better solutions.
Goldberg's metaphor: \uv{just as a child builds a castle out of simple wooden
blocks, so a GA seeks near-optimal solutions by putting building blocks
together}.
\end{defbox}

Practical consequences:
\begin{itemize}
\item \textbf{The encoding matters.} Genes that interact with each other should
  lie close together (\emph{linkage}), otherwise crossover disrupts them. Hence
  the attempts at inversion, \emph{messy GAs}, \emph{linkage learning} and
  estimation of distribution algorithms (EDA).
\item \textbf{The population size matters.} The population must contain enough
  samples of every building block (\emph{population sizing} theory).
\item \textbf{Premature convergence is harmful} --- the loss of diversity makes
  it impossible to discover the missing blocks.
\end{itemize}

\subsubsection{Implicit parallelism and the multi-armed bandit}

\begin{defbox}[Implicit parallelism]
A GA works explicitly with $n$ individuals, but implicitly it
\uv{processes} many more schemata at the same time --- between $2^m$ and
$n \cdot 2^m$. Holland estimated that the number of schemata that are processed
\emph{usefully} (they grow exponentially and are not disrupted too much) is of
the order of $n^3$.
\end{defbox}

This is often commented on as the only case in which the combinatorial explosion
works in our favour.

\paragraph{Connection with the bandit problem.} Holland's original motivation
was an \uv{adaptive plan} seeking a balance between exploration and
exploitation. Formally:
\begin{itemize}
\item \textbf{The two-armed bandit.} We have $N$ coins and a machine with two
  arms with unknown mean payoffs $\mu_1, \mu_2$ and variances. We allocate $n$
  coins to the worse arm and $N-n$ to the better one. The analytical solution
  says that the optimal strategy allocates \emph{exponentially} more trials to
  the currently winning arm: $N - n^* = O(e^{c\, n^*})$.
\item \textbf{The GA as a bandit.} The schema theorem says that a GA allocates
  exponentially more \uv{slots in the population} to the more successful
  schemata --- that is, that it solves the exploration/exploitation dilemma
  (almost) optimally.
\item It was originally believed that the SGA plays a single game with $3^m$
  arms. In fact it plays many games in parallel: the schemata compete for
  concrete fixed positions, and schemata of order $k$ play a game with $2^k$
  arms.
\end{itemize}

\subsubsection{Criticism of schema theory}

\paragraph{When a GA fails --- deceptive problems}

Consider a function in which the schemata
\[
S_1 = (111{*}\ldots{*}), \qquad S_2 = ({*}\ldots{*}11)
\]
are above average, but at the same time
\[
f(111{*}{*}{*}{*}{*}11) \;\ll\; f(000{*}{*}{*}{*}{*}00).
\]
Selection leads the GA towards the blocks of ones, even though the optimum lies
elsewhere. Such problems are called \textbf{deceptive} (Goldberg 1987): the
partial building blocks do not lead to the global optimum when put together.
Deception is sometimes proposed as a measure of the difficulty of a problem for
an EA (not everyone agrees with this).

The opposite test are the \textbf{Royal Road functions} (Mitchell, Forrest,
Holland): fitness is the sum of contributions for complete \uv{blocks} of ones
(e.g.\ eight blocks of eight bits each), so the landscape is tailored exactly to
the building blocks hypothesis and a GA ought to excel on it. The experiment
turned out the other way round: the simple GA lost to a random mutation hill
climber (RMHC). The cause is \emph{hitchhiking}: together with a good block that
has been found, selection also spreads incidental \uv{parasitic} bits at the
other positions, and these destroy the diversity needed to find the remaining
blocks. The Royal Road functions are therefore the standard empirical
counterexample to a naive reading of the BBH.

\paragraph{Static average fitness --- the static BBH}

Grefenstette (1991): people interpret the BBH as saying that a GA converges to
solutions with a good \emph{true static} average schema fitness (over the whole
space). A GA, however, estimates the fitness of a schema only \emph{from the
samples in the population}, and that estimate is biased for two reasons:

\begin{itemize}
\item \textbf{Collateral convergence.} As soon as a GA converges somewhere, it
  stops sampling the schemata uniformly. If, for instance, the schema
  $111{*}\ldots{*}$ is good, then after a few generations almost every
  individual has this prefix. But then almost every sample of the schema
  ${*}{*}{*}000\ldots$ is at the same time a sample of $111000\ldots$, and the
  fitness of the first schema is estimated in a completely distorted way.
\item \textbf{Large variance of schema fitness.} Consider $f(x) = 2$ for
  $x \in 111{*}\ldots{*}$, $f(x)=1$ for $x \in 0{*}\ldots{*}$ and $0$ otherwise.
  The static average fitness of the schema $1{*}\ldots{*}$ is $\approx 1/2$,
  whereas $F(0{*}\ldots{*}) = 1$. The GA, however, will estimate
  $F(1{*}\ldots{*}) \approx 2$, because individuals with the prefix 111 will
  prevail in the population. A high variance leads to a systematic bias and to
  the GA not sampling the schemata independently.
\end{itemize}

\paragraph{Summary of the weaknesses}

\begin{itemize}
\item The theorem is an \textbf{inequality} and only a \textbf{one-step} result.
  Iterating it over several generations requires the assumption that
  $F(S,t)/\bar f(t)$ remains constant --- which typically does not hold (both
  change dynamically).
\item Populations are \textbf{finite and small} $\Rightarrow$ sampling errors;
  the exponential growth is only an expected value.
\item The theorem accounts only for the \textbf{destructive} effect of the
  operators, not for their \emph{constructive} role (crossover can also create a
  schema).
\item The competition of schemata is \textbf{strongly dependent}, not
  independent as the arms of a bandit are.
\item The analysis is \uv{on-line}: the SGA maximises the running performance,
  so it is more suitable for adaptive problems, whereas the most common use is
  to find one best solution (an off-line criterion).
\end{itemize}

Even so, the TST remains a useful intuition: \emph{what is short, simple and
good will spread}.

\subsection{Probabilistic models of the simple GA}
\label{sec:vose}

In the 1990s (Vose, Liepins, Nix, Whitley) \emph{exact} models of the SGA were
developed, which replaced the approximate schema theory. The goal was to
describe precisely what a population looks like, to describe the transition
mapping between populations, and to study its properties and the asymptotic
behaviour of the SGA. Two approaches are distinguished: \textbf{infinite
populations} (a deterministic dynamical system) and \textbf{finite populations}
(a Markov chain).

\subsubsection{The simplified SGA}

For the analysis the SGA is simplified even further: in each step two parents
are selected, with probability $p_C$ they are crossed, \textbf{one of the
offspring is discarded}, and the remaining one is mutated and inserted into the
new population. In this way the creation of each individual of the new
population is independent.

\subsubsection{Formalisation of the population}

We identify every binary string of length $l$ with a number
$i \in \{0,1,\dots,2^l-1\}$ (e.g.\ $00000111 \equiv 7$). The population at time
$t$ is described by two vectors of length $2^l$:

\begin{defbox}[Population vector and selection vector]
\begin{itemize}
\item $p(t) = (p_0(t),\dots,p_{2^l-1}(t))$, where $p_i(t)$ is the
  \emph{relative frequency} of the string $i$ in the population. We have
  $p_i \ge 0$, $\sum_i p_i = 1$, so $p(t)$ lies in the simplex $\Lambda$.
\item $s(t) = (s_0(t),\dots,s_{2^l-1}(t))$, where $s_i(t)$ is the probability
  that selection picks the string $i$.
\end{itemize}
The vector $p$ describes the \emph{composition} of the population, the vector
$s$ the \emph{selection probabilities}.
\end{defbox}

\subsubsection{The operator $\mathcal{G}$}

The goal is to find an operator $\mathcal{G}: \Lambda \to \Lambda$ such that
\[
p(t+1) = \mathcal{G}\big(p(t)\big),
\]
and to decompose it into two parts: $\mathcal{G} = \mathcal{M} \circ
\mathcal{F}$, where $\mathcal{F}$ corresponds to selection (fitness) and
$\mathcal{M}$ to the so-called \emph{mixing} --- crossover and mutation
together. Then it suffices to iterate the operator from the initial population
and we know the complete behaviour of the SGA.

\paragraph{The part $\mathcal{F}$ --- selection.} Define the diagonal matrix
$\mathbf{F}$ with $\mathbf{F}_{ii} = f(i)$ and $\mathbf{F}_{ij} = 0$ for
$i \ne j$. Fitness-proportionate selection is then exactly
\[
s(t) \;=\; \frac{\mathbf{F}\,p(t)}{\sum_j \mathbf{F}_{jj}\,p_j(t)} .
\]
The numerator $(\mathbf{F}p)_i = f(i)\,p_i$ is the \uv{weight} of the string $i$
in the population, the denominator is a normalising constant (the average
fitness of the population). If we now drop crossover and mutation, the new
population is simply $n$ independent draws from the distribution $s(t)$, so
$\mathbb{E}[p(t+1)] = s(t)$. Substituting into the definition of selection we
obtain $\mathbb{E}[s(t+1)] \propto \mathbf{F}\,s(t)$, so iterating $\mathcal{F}$
is (up to normalisation) repeated multiplication by the matrix $\mathbf{F}$ ---
the \emph{power method}. It converges to the eigenvector belonging to the
largest diagonal value, i.e.\ to a population consisting exclusively of copies
of the best individual \emph{present in the initial population}. With finite
populations, statistical errors cause deviations from the expected value; the
larger the population, the smaller the deviation. Infinite populations are
exact, albeit impractical.

\textbf{Mini-example.} For $l = 2$ we have four strings $00,01,10,11$, i.e.\ the
indices $0,1,2,3$. Let $f = (1, 2, 3, 4)$ and let a population of eight
individuals contain $2\times 00$, $4\times 01$, $1\times 10$, $1\times 11$,
i.e.\ $p = (0.25;\,0.5;\,0.125;\,0.125)$. Then
$\mathbf{F}p = (0.25;\,1.0;\,0.375;\,0.5)$, the sum is $2.125$
(= the average fitness of the population) and
\[
s \;=\; (0.118;\; 0.471;\; 0.176;\; 0.235).
\]
The share of the best string $11$ has grown from $12.5\,\%$ to $23.5\,\%$, the
share of the worst one $00$ has dropped from $25\,\%$ to $11.8\,\%$ --- exactly
by the factor $f(i)/\bar f$, as we know from the schema theorem.

\paragraph{The part $\mathcal{M}$ --- mixing.} We introduce the function
\[
r(i,j,k) \;=\; \Pr[\text{offspring } k \text{ arises from parents } i, j],
\]
which comprises both crossover (over all cut points) and mutation. Then
\[
\mathbb{E}\big[p_k(t+1)\big] \;=\; \sum_{i}\sum_{j} s_i(t)\, s_j(t)\, r(i,j,k).
\]
This is therefore a \emph{quadratic} operator on the simplex. Deriving
$r(i,j,k)$ is technically demanding but possible. The key simplification follows
from the fact that both one-point crossover and bit mutation \uv{do not see} the
absolute values of the bits, only their agreements and disagreements. Formally:
for any $k$ we have
\[
r(i,j,k) \;=\; r(i \oplus k,\; j \oplus k,\; 0),
\]
where $\oplus$ is the bitwise XOR. It therefore suffices to know a single
\emph{mixing matrix} $\mathbf{M}$ with entries $\mathbf{M}_{ij} = r(i,j,0)$ (of
size $2^l \times 2^l$), and the rest is computed by the permutations
$\sigma_k : x \mapsto x \oplus k$:
$\mathcal{M}(s)_k = (\sigma_k s)^{\!\top}\, \mathbf{M}\, (\sigma_k s)$.
This is the main technical contribution of Vose's formalism --- the problem
shrinks from $2^{3l}$ numbers to $2^{2l}$.

\subsubsection{Results for infinite populations}

\begin{itemize}
\item The SGA is a \textbf{discrete dynamical system} on the simplex; the
  sequence $p(0), p(1), \dots$ is a trajectory of this system.
\item \textbf{Fixed points of $\mathcal{F}$}: populations formed exclusively by
  copies of the best individual from the initial population (selection
  \uv{focuses}).
\item \textbf{Fixed point of $\mathcal{M}$}: a single one --- the uniform
  distribution (mixing \uv{defocuses}, it maximises entropy).
\item The fixed points of the composed $\mathcal{G}$ are not known in general;
  finding them is the main open problem. The behaviour is a combination of
  focusing and defocusing, which manifests itself as \textbf{punctuated
  equilibria} --- long periods of metastability alternating with rapid
  transitions. This phenomenon is known from biology as well.
\end{itemize}

\subsubsection{Finite populations: a Markov chain}

\begin{defbox}[Markov chain]
A stochastic process in discrete time with states, where the probability of a
transition to the next state depends \emph{only} on the current state (not on
the history). A complete description is given by the matrix of transition
probabilities.
\end{defbox}

We model an SGA with a population of size $n$ over strings of length $l$:
\begin{itemize}
\item \textbf{States} = the individual possible populations (multisets of $n$
  strings). Their number is
  \[
  N \;=\; \binom{n + 2^l - 1}{\,2^l - 1\,},
  \]
  which is the number of multisets of size $n$ from $2^l$ types.
\item \textbf{The matrix $\mathbf{Z}$}: the columns correspond to populations,
  $\mathbf{Z}(y,i)$ = the number of occurrences of the individual $y$ in the
  population $i$. The columns of the matrix $\mathbf{Z}$ are thus the states of
  the chain.
\item \textbf{The transition matrix $\mathbf{Q}$} of size $N \times N$; it can
  be derived explicitly from $\mathcal{G}$ (the transition is a multinomial draw
  of $n$ individuals from the distribution $\mathcal{G}(p)$):
  \[
  \mathbf{Q}_{i,j} \;=\; n! \prod_{y=0}^{2^l-1}
  \frac{\big(\mathcal{G}(p_i)_y\big)^{\mathbf{Z}(y,j)}}{\mathbf{Z}(y,j)!}.
  \]
\end{itemize}

\textbf{The size problem.} Already for $n = l = 8$ the matrix $\mathbf{Q}$ has
on the order of $10^{29}$ entries. The model is therefore exact but practically
unusable for computation; its value is theoretical.

\paragraph{Qualitative consequences.}
\begin{itemize}
\item If $p_M > 0$, the chain is \textbf{irreducible and aperiodic} (from any
  population any other can be reached by mutations) $\Rightarrow$ a unique
  stationary distribution exists and the SGA visits every population (including
  the one with the global optimum) infinitely many times with probability 1.
\item Without elitism, however, it \textbf{does not converge} to the optimum ---
  the optimum appears and disappears again. An \emph{elitist} GA (one that keeps
  the best solution found) converges to the global optimum with probability 1 as
  $t \to \infty$. This is the standard convergence result; it says nothing,
  however, about the \emph{speed}.
\item There is a correspondence between the infinite- and the finite-population
  model: as $n \to \infty$ the behaviour of the finite model approaches the
  deterministic trajectory of $\mathcal{G}$.
\end{itemize}

\subsubsection{Estimation of distribution algorithms (EDA)}
\label{ssec:eda}

\noindent\mimoq{not in the EVA slides; it is the direct continuation of the Vose view}

\begin{defbox}[EDA --- the general scheme]
Vose's model shows that the SGA is in fact a machine for transforming a probability
distribution over the space of strings. \emph{Estimation of distribution algorithms} take
this literally: they throw away both crossover and mutation and maintain an
\textbf{explicit model} $p_\theta(x)$. Iteration: sample a population from $p_\theta$,
select the better part, \emph{re-learn} the model on it, repeat. The model $\theta$ is the
only information carried between generations.

By the richness of the dependences captured: \textbf{none} (UMDA, PBIL), \textbf{pairwise}
(MIMIC), \textbf{general} (BOA --- a Bayesian network), \textbf{continuous} (CMA-ES).
PBIL shifts the probability vector towards the best individual:
$p_i \gets (1-\alpha)\,p_i + \alpha\, x_i^{\mathrm{best}}$.
\end{defbox}

A univariate model fails as soon as there is \emph{epistasis} between the genes --- hence
BOA, which learns the structure of the dependences. EDAs thus solve the problem for the sake
of which Holland proposed inversion: where the cooperating genes should lie is something the
algorithm learns by itself.

\subsection{Exam summary}

\begin{itemize}
\item Schema = a word over $\{0,1,*\}$; $o(S)$ = the number of fixed positions;
  $d(S)$ = the distance between the first and the last fixed position;
  $F(S)$ = the average fitness of the representatives \emph{in the given
  population}. There are $3^m$ schemata in total, and an individual is a
  representative of $2^m$ of them.
\item TST: $\mathbb{E}[C(S,t+1)] \ge C(S,t)\frac{F(S)}{\bar f}
  [1 - p_C \frac{d(S)}{m-1} - p_M o(S)]$; the derivation in three steps
  (selection $\to$ exponential growth, crossover $\to$ $d(S)/(m-1)$,
  mutation $\to$ $(1-p_M)^{o(S)}$).
\item BBH: a GA assembles short, low-order, above-average schemata (building
  blocks). Consequences: the encoding matters, the population size matters,
  premature convergence is harmful.
\item Implicit parallelism ($\sim n^3$ usefully processed schemata), the analogy
  with the $k$-armed bandit (exponential allocation of trials to the winner).
\item Criticism: an inequality for a single step, finite populations, deceptive
  problems, collateral convergence, large variance of fitness, disregard of the
  constructive effects, dependence of the schemata.
\item Vose's model: strings are identified with the numbers $0..2^l-1$; the
  population is described by the vector of relative frequencies $p(t)$ in the
  simplex and by the vector of selection probabilities $s(t)$.
\item We look for $\mathcal{G} = \mathcal{M} \circ \mathcal{F}$ with
  $p(t+1)=\mathcal{G}(p(t))$. $\mathcal{F}$ (selection) is given by the diagonal
  fitness matrix: $s = \mathbf{F}p / \sum_j \mathbf{F}_{jj}p_j$;
  $\mathcal{M}$ (crossover + mutation) is a quadratic operator
  $\mathbb{E}[p_k'] = \sum_{i,j} s_i s_j\, r(i,j,k)$. Thanks to the symmetry
  $r(i,j,k) = r(i\oplus k, j \oplus k, 0)$ a single mixing matrix suffices.
\item Fixed points: $\mathcal{F}$ --- a population of copies of the best
  individual; $\mathcal{M}$ --- the uniform distribution. The combination leads
  to punctuated equilibria. The fixed points of $\mathcal{G}$ are not known.
\item Finite populations: a Markov chain with states = populations,
  $N = \binom{n+2^l-1}{2^l-1}$ states; the matrix $\mathbf{Q}$ is exact but
  astronomically large. Irreducibility for $p_M>0$; an elitist GA converges to
  the optimum with probability 1.
\item Infinite populations = a deterministic model (exact, impractical), finite
  ones = a stochastic model (realistic, intractable). The limit transition
  connects them.
\item \textbf{EDAs} take the view \uv{a GA = a transformation of a
  distribution} literally: instead of operators a model $p_\theta$ is learned
  (UMDA/PBIL --- independent marginals, MIMIC/BMDA --- pairwise dependences,
  BOA --- a Bayesian network, CMA-ES --- a Gaussian). The cycle: sample $\to$
  select the better ones $\to$ re-learn the model. PBIL:
  $p_i \gets (1-\alpha)p_i + \alpha x_i^{\mathrm{best}}$.
\end{itemize}

%=====================================================================
\section[Evolution strategies, differential evolution, coevolution, open-ended evolution]%
{Evolution strategies, differential evolution,\protect\\
coevolution, open-ended evolution}
\label{sec:es}
%=====================================================================

The chapter joins four themes that have in common that they go \emph{beyond} the classical
genetic algorithm: two methods for continuous optimisation (ES and DE) and two themes in
which the fitness itself changes --- coevolution (fitness depends on the other individuals)
and open-ended evolution (fitness is not fixed at all).

\subsection{Evolution strategies}
\label{ssec:es}
\subsubsection{Motivation and basic features}

Evolution strategies (ES) were proposed by Rechenberg and Schwefel in Berlin in
the 1960s for the optimisation of real-valued parameters in difficult
engineering problems (the shape of a nozzle, the bend of a pipe).
Characteristics:

\begin{itemize}
\item An individual is a vector of real numbers; \textbf{mutation is the main
  operator}.
\item Besides the \emph{genetic parameters} (the solution itself) an individual
  also contains \emph{strategy parameters}, also called \emph{endogenous
  parameters}, which control the evolution itself --- typically the standard
  deviations of the mutations. Hence the slogan \uv{the evolution of
  evolution}.
\item Parent selection is \textbf{random} (independent of fitness);
  environmental selection is \textbf{deterministic} --- the $\mu$ best are
  selected.
\item The most modern and most successful representative is CMA-ES.
\end{itemize}

\begin{defbox}[ES notation]
$\mu$ --- the population size (the number of parents), $\lambda$ --- the number
of offspring created, $\rho$ --- the number of parents taking part in one
offspring.
\begin{itemize}
\item $(\mu + \lambda)$-ES: the new population is selected from \emph{both the
  $\mu$ parents and the $\lambda$ offspring} (elitist).
\item $(\mu, \lambda)$-ES: the new population is selected \emph{from the
  $\lambda$ offspring only}; the parents always die out (non-elitist),
  necessarily $\lambda > \mu$.
\item $(\mu/\rho \,\overset{+}{,}\, \lambda)$-ES: an explicit notation with the
  recombination of $\rho$ parents.
\end{itemize}
\end{defbox}

\begin{center}
\begin{tabular}{@{}p{4cm}p{5.5cm}p{5.5cm}@{}}
\toprule
 & $(\mu,\lambda)$ & $(\mu+\lambda)$ \\
\midrule
pool for selection & offspring only & parents + offspring \\
elitism & no (fitness may decrease) & yes \\
escape from local optima & better & worse \\
convergence & slower & faster \\
recommendation & continuous domains, noisy or dynamic fitness & discrete
  domains, expensive fitness \\
self-adaptation of $\sigma$ & works well (a bad $\sigma$ dies out) & may get
  stuck (an old parent with a small $\sigma$ survives) \\
ratio & typically $\lambda/\mu \approx 7$ & even $\mu \ge \lambda$;
  $\lambda = 1$ = steady-state ES \\
\bottomrule
\end{tabular}
\end{center}

\subsubsection{$(1+1)$-ES and the 1/5 rule}

The simplest ES: the population is a single individual, a single offspring is
created by Gaussian mutation and is accepted only if it is better.

\begin{algorithm}[!ht]
\caption{$(1+1)$-ES with the 1/5 rule (minimisation of $f:\mathbb{R}^n \to \mathbb{R}$)}
\begin{algorithmic}[1]
\State initialise $\vec x(0)$ at random, $\sigma > 0$, $t \gets 0$
\Repeat
  \State $\vec y(t) \gets \vec x(t) + \sigma \cdot N(\vec 0, \mathbf{I})$
  \If{$f(\vec y(t)) < f(\vec x(t))$} $\vec x(t+1) \gets \vec y(t)$
  \Else\ $\vec x(t+1) \gets \vec x(t)$
  \EndIf
  \State track $p$ = the proportion of successful mutations over the last $k$ iterations
  \State every $k$ iterations adjust $\sigma$:
  \State \quad \textbf{if} $p > 1/5$ \textbf{then} $\sigma \gets \sigma / c$
    \Comment{going well: take larger steps, more exploration}
  \State \quad \textbf{if} $p < 1/5$ \textbf{then} $\sigma \gets \sigma \cdot c$
    \Comment{going badly: shorten the steps, more exploitation}
  \State \quad \textbf{if} $p = 1/5$ \textbf{then} leave $\sigma$ unchanged
    \Comment{typically $0.8 \le c \le 1$}
  \State $t \gets t+1$
\Until{$\vec x(t)$ is good enough}
\end{algorithmic}
\end{algorithm}

\begin{defbox}[The 1/5 success rule]
A heuristic derived by Rechenberg from the analysis of two model functions (the
sphere and the corridor): \emph{the optimal rate of successful mutations is
approximately $1/5$.} If the success rate is higher, the step is too small (we
are being too cautious) and $\sigma$ is increased; if it is lower, the step is
too large and $\sigma$ is decreased. This is \emph{adaptive} (not
self-adaptive) parameter control --- it uses feedback from the run, but the
parameter is not part of the genome.
\end{defbox}

\subsubsection{Self-adaptation of the strategy parameters}

\begin{defbox}[Self-adaptation]
The strategy parameters (the deviations $\sigma$, possibly rotation angles) are
\textbf{part of the genome} and are subject to mutation and selection together
with the solution. The selection is indirect: individuals with a suitable
$\sigma$ produce better offspring, and thereby their $\sigma$ spreads through
the population. The order is essential: \textbf{first the strategy parameters
are mutated, and only then are the genetic parameters mutated using them} ---
otherwise the old $\sigma$ would be the one being evaluated.
\end{defbox}

An individual has the form $C = [\vec x, \vec \sigma]$, where
$\vec x \in \mathbb{R}^n$ and $\vec\sigma$ has $k$ components:

\begin{center}
\begin{tabular}{@{}clp{8.2cm}@{}}
\toprule
$k$ & name & geometry of the mutation \\
\midrule
$1$ & one global $\sigma$ & isotropic --- the density contours are spheres \\
$n$ & uncorrelated mutations & an ellipsoid with axes parallel to the
  coordinates \\
$n(n+1)/2$ & correlated mutations & a general ellipsoid ($n$ deviations $+$
  $n(n-1)/2$ rotation angles $\alpha_{ij}$), corresponding to a full covariance
  matrix \\
\bottomrule
\end{tabular}
\end{center}

The covariance matrix $\mathbf{C}$ is a product with a rotation matrix:
$c_{ii} = \sigma_i^2$,
$c_{ij} = \tfrac12(\sigma_i^2 - \sigma_j^2)\tan(2\alpha_{ij})$.

\paragraph{The standard mutation rules (Schwefel).}
\[
\sigma_i' = \sigma_i \cdot \exp\big(\tau' N(0,1) + \tau N_i(0,1)\big),
\qquad
x_i' = x_i + \sigma_i' \cdot N_i(0,1),
\]
where $\tau' \propto 1/\sqrt{2n}$ (a common perturbation for the whole
individual) and $\tau \propto 1/\sqrt{2\sqrt{n}}$ (an individual perturbation of
the component). The lognormal form guarantees the positivity of $\sigma$ and the
symmetry of increases and decreases. A lower bound $\sigma_{\min}$ is introduced
so that $\sigma$ does not degenerate to zero. The angles are mutated additively:
$\alpha_{ij}' = \alpha_{ij} + \beta \cdot N(0,1)$.

\subsubsection{Recombination in ES}

\emph{One} offspring is created from $\rho$ parents. According to the number of
parents we distinguish \textbf{local} ($\rho = 2$) and \textbf{global}
($\rho = \mu$) recombination; according to the way of combining:
\begin{itemize}
\item \textbf{Discrete (dominant):} the value at a given position is chosen at
  random from one of the parents.
\item \textbf{Intermediate (arithmetic):} the value is the average of the
  parents' values at the given position.
\end{itemize}
The usual combination: discrete recombination for the genetic parameters (it
preserves diversity) and intermediate recombination for the strategy parameters
(it stabilises the adaptation of $\sigma$).

\begin{algorithm}[!ht]
\caption{The general ES cycle}
\begin{algorithmic}[1]
\State $t \gets 0$; initialise at random a population $P_t$ of $\mu$ individuals $[\vec x, \vec\sigma]$
\State evaluate the individuals in $P_t$
\While{the termination condition is not met}
  \State $C_{t+1} \gets \emptyset$
  \For{$\lambda$ times}
    \State select $\rho$ parents at random from $P_t$
    \State recombine them into a single offspring $[\vec x, \vec\sigma]$
    \State mutate the strategy parameters $\vec\sigma$
    \State mutate the genetic parameters $\vec x$ using the new $\vec\sigma$
    \State evaluate the offspring and insert it into $C_{t+1}$
  \EndFor
  \State $(\mu,\lambda)$: $P_{t+1} \gets$ the $\mu$ best of $C_{t+1}$
  \State $(\mu+\lambda)$: $P_{t+1} \gets$ the $\mu$ best of $P_t \cup C_{t+1}$
  \State $t \gets t+1$
\EndWhile
\end{algorithmic}
\end{algorithm}

\subsubsection{CMA-ES}

\noindent\mimoq{the slides give it one line as \uv{today's most successful (and complex)} ES}

\begin{defbox}[CMA-ES --- the principle]
\emph{Covariance Matrix Adaptation ES} (Hansen, Ostermeier) is nowadays the reference
algorithm for continuous black-box optimisation. Instead of self-adapting the individual
$\sigma$ values in the genome it maintains \textbf{a single distribution for the whole
population}, $\mathcal{N}(\vec m, \sigma^2 \mathbf{C})$, and learns its covariance matrix
from the history of successful steps.

The iteration: sample $\lambda$ points, select the $\mu$ best \emph{by rank}, move the centre
$\vec m$ to their weighted average, and update $\mathbf{C}$ and the step length $\sigma$
according to the \emph{evolution paths} --- exponentially smoothed records of the direction
in which the centre has been moving. If the steps go in a consistent direction, $\sigma$ is
increased; if they oppose each other, it is decreased.
\end{defbox}

The result is an algorithm that adapts automatically to the anisotropy and the correlations
of the problem (in essence it learns an approximation of the inverse Hessian without a
gradient), is invariant with respect to monotone transformations of fitness, and has very few
tunable parameters.

\subsection{Differential evolution}
\label{sec:de}
\subsubsection{Motivation}

Differential evolution (DE; Storn and Price, 1997) is a simple and surprisingly
powerful algorithm for the black-box optimisation of functions
$f: \mathbb{R}^D \to \mathbb{R}$. The key idea is \textbf{geometric}: neither
the size nor the direction of the mutation step is estimated from a parameter
$\sigma$; both are taken over from the \emph{differences of pairs of individuals
in the population}. The population thus encodes the scale and the orientation of
the problem by itself: at the beginning it is scattered (large steps), in a
converging state the differences are small (fine tuning). It is a kind of
implicit self-adaptation for free.

The applications range from machine learning to the planning of optimal
trajectories of space probes (the European Space Agency).

\subsubsection{The DE/rand/1/bin algorithm}

The population consists of $N$ vectors
$\vec x_{1},\dots,\vec x_{N} \in \mathbb{R}^D$. For \textbf{every} individual
$\vec x_i$ (parent selection is thus trivial --- every individual becomes a
parent in turn) the following is performed:

\begin{enumerate}
\item \textbf{Mutation.} Select three mutually distinct individuals
  $\vec x_{a}, \vec x_{b}, \vec x_{c}$, all different from $\vec x_i$, and
  create the \emph{donor vector}
  \[
  \vec v_i \;=\; \vec x_{a} \;+\; F\,(\vec x_{b} - \vec x_{c}),
  \qquad F \in [0,2],\ \text{typically } F \approx 0.8 .
  \]
\item \textbf{Crossover (binomial, uniform \uv{with a safeguard}).}
  Create the \emph{trial vector} $\vec u_i$:
  \[
  u_{j,i} =
  \begin{cases}
  v_{j,i} & \text{if } \mathrm{rand}_j \le C \text{ or } j = I_{\mathrm{rand}},\\[2pt]
  x_{j,i} & \text{otherwise},
  \end{cases}
  \]
  where $\mathrm{rand}_j \sim U(0,1)$, $C \in [0,1]$ is the crossover threshold
  and $I_{\mathrm{rand}}$ is a random index from $\{1,\dots,D\}$. The condition
  $j = I_{\mathrm{rand}}$ is the \uv{safeguard}: it guarantees that at least one
  component comes from the donor, so the offspring is never an identical copy of
  the parent.
\item \textbf{Environmental selection (a 1:1 duel).}
  \[
  \vec x_i(t+1) =
  \begin{cases}
  \vec u_i & \text{if } f(\vec u_i) \le f(\vec x_i(t)),\\
  \vec x_i(t) & \text{otherwise.}
  \end{cases}
  \]
  Selection is therefore elitist at the level of the individual --- the
  population never gets worse.
\end{enumerate}

\begin{algorithm}[!ht]
\caption{DE/rand/1/bin}
\begin{algorithmic}[1]
\Require population size $N$, scale factor $F$, crossover threshold $C$
\State initialise the population $\{\vec x_i\}_{i=1}^{N}$ at random within the bounds of the problem
\While{the termination condition is not met}
  \For{$i \gets 1$ \textbf{to} $N$}
    \State select distinct random indices $a,b,c \ne i$
    \State $\vec v \gets \vec x_a + F(\vec x_b - \vec x_c)$
    \State $I_{\mathrm{rand}} \gets$ a random index from $\{1,\dots,D\}$
    \For{$j \gets 1$ \textbf{to} $D$}
      \State $u_j \gets v_j$ if $\mathrm{rand}() \le C$ or $j = I_{\mathrm{rand}}$, otherwise $u_j \gets x_{j,i}$
    \EndFor
    \If{$f(\vec u) \le f(\vec x_i)$} $\vec x_i' \gets \vec u$ \Else\ $\vec x_i' \gets \vec x_i$ \EndIf
  \EndFor
  \State $\{\vec x_i\} \gets \{\vec x_i'\}$
\EndWhile
\end{algorithmic}
\end{algorithm}

\subsubsection{A numerical example of one step}

Let $D = 3$ and let us minimise $f(\vec x) = \sum_j x_j^2$ with the parameters
$F = 0.8$, $C = 0.9$. The current individual is
$\vec x_i = (4;\, -3;\, 2)$, $f(\vec x_i) = 16+9+4 = 29$.
The randomly selected individuals are:
\[
\vec x_a = (2;\, -1;\, 0.5), \quad
\vec x_b = (1;\, 0;\, -1), \quad
\vec x_c = (-1;\, 2;\, 1).
\]

\textbf{Mutation:}
\[
\vec x_b - \vec x_c = (2;\, -2;\, -2), \qquad
\vec v = (2;\,-1;\,0.5) + 0.8\,(2;\,-2;\,-2) = (3.6;\; -2.6;\; -1.1).
\]

\textbf{Crossover:} let $I_{\mathrm{rand}} = 1$ and let the random numbers be
$\mathrm{rand} = (0.40;\ 0.95;\ 0.20)$. Then
\begin{itemize}
\item $j=1$: $0.40 \le 0.9$ (and moreover $j = I_{\mathrm{rand}}$) $\Rightarrow u_1 = v_1 = 3.6$;
\item $j=2$: $0.95 > 0.9$ and $j \ne I_{\mathrm{rand}}$ $\Rightarrow u_2 = x_{2,i} = -3$;
\item $j=3$: $0.20 \le 0.9$ $\Rightarrow u_3 = v_3 = -1.1$.
\end{itemize}
The trial vector is $\vec u = (3.6;\, -3;\, -1.1)$.

\textbf{Selection:} $f(\vec u) = 12.96 + 9 + 1.21 = 23.17 < 29 = f(\vec x_i)$,
so $\vec u$ replaces $\vec x_i$ in the new population.

\subsubsection{Mutation variants}

The notation \texttt{DE/$x$/$y$/$z$}: $x$ = how the base vector is chosen,
$y$ = the number of difference vectors, $z$ = the type of crossover
(\texttt{bin} = binomial uniform, \texttt{exp} = exponential, a contiguous
segment).

\begin{center}
\begin{tabular}{@{}ll@{}}
\toprule
\textbf{variant} & \textbf{donor $\vec v$} \\
\midrule
DE/rand/1 & $\vec x_a + F(\vec x_b - \vec x_c)$ \\
DE/rand/2 & $\vec x_a + F(\vec x_b - \vec x_c + \vec x_d - \vec x_e)$ \\
DE/best/1 & $\vec x_{\mathrm{best}} + F(\vec x_b - \vec x_c)$ \\
DE/best/2 & $\vec x_{\mathrm{best}} + F(\vec x_b - \vec x_c + \vec x_d - \vec x_e)$ \\
DE/current-to-best/1 & $\vec x_i + F(\vec x_{\mathrm{best}} - \vec x_i) + F(\vec x_b - \vec x_c)$ \\
\bottomrule
\end{tabular}
\end{center}

The variants with \texttt{best} converge faster, but they get stuck in a local
optimum more easily; \texttt{rand} is more robust. The modern variants (jDE,
SaDE, JADE, SHADE, L-SHADE) adapt $F$ and $C$ during the run and are among the
top performers in the CEC competitions.

\subsubsection{Comparison of DE with ES and GA}

\begin{center}
\begin{tabular}{@{}p{3.2cm}p{3.7cm}p{3.7cm}p{3.7cm}@{}}
\toprule
 & \textbf{GA (real-valued)} & \textbf{ES} & \textbf{DE} \\
\midrule
source of the mutation step & fixed $\sigma$ & evolved $\sigma$ & differences of
  individuals \\
parent selection & according to fitness & random & every individual once \\
env.\ selection & generational / elitism & the $\mu$ best & duel parent vs.\
  offspring \\
main operator & crossover & mutation & differential mutation \\
adaptation of the scale & none & explicit (in the genome) & implicit (from the
  population) \\
\bottomrule
\end{tabular}
\end{center}

\subsection{Coevolution}
\label{sec:koevoluce}
\subsubsection{Contextual fitness}

\begin{defbox}[Coevolution]
An evolutionary algorithm in which the fitness of an individual is \textbf{not
absolute} but depends on the other individuals --- on its own population or on
another population evolving alongside it. Fitness is therefore
\emph{contextual}, and the fitness landscape is \textbf{dynamic}: it changes as
the opponents or the partners change.
\end{defbox}

Motivation: for many problems we cannot write down an absolute objective
function (\uv{how good is a strategy for draughts?}), but we can \emph{compare}
two individuals (by playing a game). In addition, coevolution promises
\emph{evolutionary arms races} --- the opponents push each other towards ever
better solutions. Types of coevolution:

\begin{center}
\begin{tabular}{@{}p{3cm}p{6cm}p{6cm}@{}}
\toprule
 & \textbf{competitive} & \textbf{cooperative} \\
\midrule
relationship & the fitness of one grows at the expense of the other &
  the individuals complement each other into a common solution \\
example & predator--prey, host--parasite, player--opponent, sorting networks
  vs.\ sequences of numbers & decomposition of a problem into subproblems,
  modules and blueprints (CoDeepNEAT), main programs and ADFs \\
fitness & the outcome of the duel (the tournament) & the quality of the
  composite solution the individual took part in \\
risk & cycling, disengagement, overspecialisation & credit assignment, lazy
  partners \\
\bottomrule
\end{tabular}
\end{center}

\subsubsection{Competitive coevolution}

\paragraph{The classical example: Hillis's sorting networks (1990).} A
population of sorting networks coevolves against a population of test sequences.
The fitness of a network = the proportion of correctly sorted tests; the fitness
of a test = how many networks failed on it. The parasitic tests specialise in
the weaknesses of the networks, and the networks have to improve. In this way
Hillis found a network for 16 elements with \textbf{61} comparators.
\emph{Beware of a frequent mistake:} the best known \emph{human} design (Green,
1969) has 60 comparators, so coevolution did not beat it. What matters is the
comparison within the experiment --- the same genetic algorithm \emph{without}
coevolving parasites achieved only 65 comparators. Coevolution therefore
demonstrably improved the optimisation ability of the EA and came within a
single comparator of the human result.

\paragraph{Predator--prey} in evolutionary robotics: hunter robots and fugitive
robots perfect each other. \textbf{Host--parasite}: models in artificial life
(Tierra), software testing.

\paragraph{Tournament evaluation.} Fitness is obtained from duels: round robin
(expensive), a random sample of opponents (cheap, noisy), a \emph{hall of fame}
(a duel against an archive of the historically best individuals --- it prevents
forgetting), or a \emph{tournament with shared fitness} (\emph{competitive
fitness sharing}: defeating an opponent whom few can defeat brings a higher
reward).

\subsubsection{Pathologies of coevolution}

\begin{itemize}
\item \textbf{Cycling} (intransitivity): the populations go round in a circle as
  in rock--paper--scissors; the apparent progress is not real. The dynamics of
  coevolution can even be deterministically chaotic.
\item \textbf{Disengagement}: one population completely outgrows the other, all
  duels end the same way, fitness loses its discriminating power and selection
  stops working (the gradient disappears).
\item \textbf{The Red Queen effect}: everyone improves, but the relative fitness
  stays constant --- from the measured values we cannot tell whether the
  progress is real. The remedy: measuring against a fixed external set of
  opponents (a \emph{master tournament}).
\item \textbf{Mediocre stable states}: the populations settle into a mutual
  equilibrium at a low level.
\item \textbf{Overspecialisation} (\emph{focusing}): an individual optimises
  itself against the current opponent and loses generality.
\item \textbf{Forgetting}: the ability to defeat old opponents is lost. The
  remedy: archives (hall of fame, Pareto-coevolutionary archives, \emph{Nash
  memory}).
\end{itemize}

\subsubsection{Cooperative coevolution}

\begin{defbox}[CCGA (cooperative coevolutionary GA, Potter \& De Jong)]
The problem is split into $k$ components (e.g.\ coordinates, modules, rules).
A \textbf{separate population} is maintained for each component. The fitness of
an individual from population $i$ is determined by \emph{combining} it with
representatives of the other populations (typically with their currently best
individuals, possibly also with random ones) and evaluating the resulting
composite solution.
\end{defbox}

Advantages: the decomposition of a large space into smaller ones, natural
parallelisation, suitability for modular problems. Problems: \emph{credit
assignment} --- how to divide the quality of the composite solution among the
components; \emph{interdependence} (strongly interacting components cannot be
separated well); \emph{relative overspecialisation} (a population adapts to
particular partners).

Examples in practice: the coevolution of main programs and ADFs in GP,
CoDeepNEAT (a population of modules and a population of \uv{blueprints},
section~\ref{sec:neuro}), meta-Lamarckian memetic algorithms (a population of
memes and a population of solutions, chapter~\ref{sec:lokalni}).

\subsubsection{The evolution of cooperation and the prisoner's dilemma}

Coevolution is closely related to a basic question of evolutionary biology:
\emph{how can altruism arise} when by definition it lowers the fitness of the
individual? Darwinism is competitive in its essence (\uv{survival of the
fittest}), and yet there is a great deal of cooperation both in nature and in
society. The theories:
\begin{itemize}
\item \textbf{Group selection} --- evolution acts on groups as well (Darwin);
  the problem: how to explain cheaters inside the group.
\item \textbf{Kin selection} --- the preservation of almost identical genes in
  relatives; the problem: altruism towards strangers and other species.
\item \textbf{The selfish gene} (Dawkins) --- the unit of evolution is the gene,
  not the individual. (Wilson: \uv{an organism is only DNA's way of making more
  DNA}.)
\item \textbf{Reciprocal altruism} (Trivers, 1971) --- mutual benefit, the
  \emph{shadow of the future}; the formal model = the iterated prisoner's
  dilemma.
\end{itemize}

\begin{defbox}[The prisoner's dilemma]
A two-player game with the actions $C$ (cooperate) and $D$ (defect) and the
payoffs
\[
\begin{array}{c|cc}
 & C & D\\\hline
C & R/R & S/T\\
D & T/S & P/P
\end{array}
\qquad
T > R > P > S \quad (\text{and for the iterated version } 2R > T+S).
\]
Typically $R=-1$, $T=0$, $P=-2$, $S=-3$. Because $T > R$ and $P > S$, defection
is the \textbf{dominant strategy} of both players; $DD$ is a \textbf{Nash
equilibrium}, but as the only outcome it \textbf{is not Pareto-optimal} ---
$CC$ maximises the joint gain.
\end{defbox}

Rational players thus arrive at an outcome that is worse for both of them than
the cooperation they could have achieved. The same mechanism in its $n$-player
form is called the \emph{tragedy of the commons}: it individually pays everyone
to draw a little more from a shared resource, and the resource therefore
perishes. The question \uv{what is actually rational (and do people behave that
way?)} is at the heart of the whole discussion around the prisoner's dilemma.

\begin{defbox}[Dominant strategy, Nash equilibrium, Pareto optimality]
A strategy $s$ is \emph{dominant} for player $i$ if it gives the same or a
better outcome than any other strategy of $i$ against \emph{all} strategies of
the opponent. A pair $(s_i, s_j)$ is a \emph{Nash equilibrium} if the strategies
are mutual best responses (no player has a reason to change strategy
unilaterally). A solution is \emph{Pareto-optimal} if the outcome of one player
cannot be improved without worsening the outcome of another.
Nash's theorem: every game with finitely many strategies has an equilibrium in
\emph{mixed} strategies (a random choice among the pure ones) --- e.g.\
rock--paper--scissors.
\end{defbox}

\paragraph{The iterated version and Axelrod's tournaments.} If the game is
played repeatedly and the players remember the history, cooperation may pay off
(the \emph{shadow of the future}). If the exact number of rounds $N$ is known in
advance, backward induction shows that it is optimal to keep defecting --- in
practice, however, the players do not know the end. Axelrod (1980) organised
tournaments of strategies:
\begin{itemize}
\item 1st tournament: 14 strategies + RANDOM, 200 games, round robin (including
  against itself), 5 repetitions. \textbf{TFT} (\emph{tit for tat}) won: start
  by cooperating, then repeat the opponent's last move.
\item 2nd tournament: 62 strategies, everyone knew the results of the first ---
  TFT won again.
\item 3rd, the \uv{ecological} tournament: a simulation as in a GA --- the
  number of individuals of a strategy in the next generation is proportional to
  the number of wins, 1000 generations. TFT prevailed once more.
\end{itemize}

\paragraph{Four properties of successful strategies:}
\emph{niceness} (do not defect first),
\emph{provocability} (retaliation --- punish a defection immediately),
\emph{forgiveness} (but calm down quickly),
\emph{non-enviousness} (do not strive to beat a \emph{particular} opponent; TFT
never scored more points than its opponent, and yet it won the tournament).
As a fifth property Axelrod lists \emph{clarity} (be simple, so that others are
able to establish cooperation with you at all).
No strategy wins against everybody; one has to be successful against very
diverse opponents (ALL-D, TFTT, RANDOM, TRIGGER) and to play well against
oneself.

\paragraph{Lessons and later developments.} In environments where the payoffs
favour cooperation and where the probability of repetition is high, cooperation
usually evolves --- and neither rationality nor consciousness is needed for
that, higher fitness suffices. In a noisy environment (a move executed
incorrectly) the more successful strategy is \textbf{Pavlov}
(\emph{win-stay, lose-shift}): if I received $R$ or $P$, repeat the move ($C$);
if I received $T$ or $S$, change it ($D$).
The tournament repeated after 20 years was won by a \emph{team} of strategies
from the University of Southampton: the first ten moves or so served to
recognise a teammate, after which all the members of the team selflessly helped
one \uv{parent} strategy to obtain a high score.

\subsubsection{Diversity, species and niches}

A small selection pressure is enough for the diversity of a population to
disappear. (A similar line of thought is developed by Flegr's \emph{frozen
evolution}: a genetically variable species behaves flexibly towards selection
only for a short time, after which it \uv{freezes} and stops responding to
changes of the environment.) Mechanisms for maintaining diversity (useful for
coevolution as well as for dynamic fitness and multi-objective optimisation):
\begin{itemize}
\item \textbf{Fitness sharing:} fitness is divided by the number of similar
  individuals,
  \[
  f'(i) = f(i)\Big/\sum_j sh\big(d(i,j)\big),
  \qquad
  sh(d) = 1-(d/\sigma_{\mathrm{share}})^{\alpha}
  \]
  for $d < \sigma_{\mathrm{share}}$ and 0 otherwise. It creates pressure
  towards occupying different niches.
\item \textbf{Crowding:} a new individual replaces the \emph{most similar}
  individual (deterministic crowding, RTS).
\item \textbf{Non-random mating:} crossover only between similar individuals
  (\emph{speciation}), or according to a topology --- the individuals live on a
  2D/3D grid and mate only with their neighbours.
\item \textbf{The island model:} several subpopulations with occasional
  migration (in more detail below).
\item \textbf{Explicit species} by means of tag bits, or according to the
  structural similarity of the genomes (NEAT, section~\ref{sec:neuro}); mating
  is then allowed only within a species.
\end{itemize}
The problem is always the definition of similarity (Hamming distance, distance
in the phenotype space, the number of matching innovation numbers) and the
choice of the niche threshold.

\paragraph{Parallel models of EAs.} The structure of the population is closely
related to parallelisation; three models are distinguished:
\begin{itemize}
\item \textbf{Master--slave} (global parallelisation): a single population, and
  only the \emph{fitness evaluation} is distributed among the nodes. The
  algorithm is mathematically identical with the sequential one; it pays off
  when fitness is expensive.
\item \textbf{The island / coarse-grained model}: several independent
  subpopulations between which a few individuals occasionally \emph{migrate}.
  Parameters: the topology of the islands, the \emph{migration interval} and the
  \emph{migration rate}, the choice of the migrants and of those they replace.
  Rare migration maintains diversity (the islands converge to different places);
  if it is too frequent, the model degenerates into a single population.
\item \textbf{The cellular / fine-grained (diffusion) model}: the individuals sit
  on a 2D/3D grid and mate only with their neighbours. A good solution spreads
  through the population like a wave, which strongly slows down premature
  convergence; this is an implicit form of non-random mating.
\end{itemize}
The island and the cellular model are therefore not merely an implementation
trick --- they change the very dynamics of the algorithm and belong among the
standard tools for maintaining diversity.

\subsubsection{Dynamic fitness landscapes}

\begin{defbox}[Fitness landscape]
Introduced by Sewall Wright (1932). A graphical representation of the dependence
of fitness on the genotype (or on the allele frequencies or on a phenotypic
trait). A theoretical tool for reasoning about EAs; in higher dimensions,
however, it is not intuitive.
\end{defbox}

Fitness often changes in time: a robot in a room where somebody turns the light
on; it starts raining; a stock market crash; tournaments of agents
(coevolution). Classical GAs perform poorly, because they are \uv{overfitted} to
static problems --- they converge quickly and lose the diversity they would need
in order to react to the change.
De Jong's comparison: a $(1+10)$-ES reacts slowly to a sudden change, because
self-adaptation has shrunk its step; a generational GA with roulette wheel
selection and without elitism keeps more diversity and recovers faster.

Modifications of EAs for dynamic landscapes:
\begin{itemize}
\item \textbf{Diploid representation} (Goldberg, Smith) --- two sets of
  chromosomes with dominance act as a long-term memory when the fitness
  oscillates.
\item \textbf{Hypermutation} (Cobb, Grefenstette) --- if the average fitness of
  the population drops (the environment has probably changed), the mutation
  probability is raised considerably.
\item \textbf{Random immigrants} --- a batch of new random individuals is
  inserted into the population.
\item \textbf{Maintaining diversity} --- a lower selection pressure, crowding,
  niching, subpopulations, the comma ES $(\mu,\lambda)$ (the parents do not
  survive, so the population detaches itself from an obsolete optimum more
  easily).
\end{itemize}
Types of change: small and smooth (wear of the hardware), morphological (hills
appear and disappear), cyclic (the seasons), catastrophic and discontinuous.
If the fitness changes too fast, no solution is permanently good
(cf.\ No Free Lunch).

\subsection{Open-ended evolution}
\label{ssec:otevrena}

\noindent\mimoq{the whole section except interactive evolution} The topic names open-ended
evolution explicitly, but neither the EVA~I nor the EVA~II slides contain anything on it,
and \emph{Evolutionary Robotics} mentions it in a single sentence (\uv{the evolution of
robots is a continuous process with an open end}). It is therefore treated briefly, at a
level that will do at the examination. The exception is \emph{interactive evolution} at the
end of the section --- that one is in the Evolutionary Robotics slides (subjective fitness,
Biomorphs).

\begin{defbox}[Open-ended evolution]
An evolutionary process that \textbf{has no fixed goal} and that keeps producing new
complexity and new \uv{kinds} of behaviour without converging. The model is the terrestrial
biosphere: it has no objective function, and yet it permanently generates novelty.
\end{defbox}

\paragraph{The hallmarks of open-endedness.} In order to be able to decide about a system at
all, \emph{hallmarks} are formulated: \textbf{the sustained production of novelty} (new
behaviour appears even after an arbitrarily long run); \textbf{the unbounded growth of
complexity} of the individuals and of their relationships; \textbf{the emergence of new
classes of entities} and levels of organisation (in biology the \emph{major transitions}:
replicator $\to$ cell $\to$ multicellularity $\to$ society); \textbf{the evolution of
evolvability} --- the very way in which the system generates novelty changes as well.

\paragraph{The systems and why they get stuck.} \emph{Tierra} (Ray) and \emph{Avida} ---
self-copying machine-code programs in which parasites and hyperparasites emerged;
\emph{Polyworld} --- agents with neural networks in a 2D world; \emph{ECHO} (Holland) ---
agents with \uv{tags} and resources. Practically every one of them reaches, after some time,
a state in which nothing qualitatively new appears any more: the environment is finite, the
space of possible \uv{inventions} is exhausted and the dynamics freezes. The hypotheses as to
why: the lack of a sufficiently rich and self-extending environment, the lack of a genuinely
open representation, and the lack of a mechanism that would continually raise the demands
placed on the individuals. Achieving genuinely open-ended evolution remains an open research
problem.

\paragraph{Generative (indirect) representations.} Related to open-ended evolution is the
question of how to encode unboundedly growing complexity at all. The answer are encodings in
which the genotype is not a description of the result but \textbf{a recipe for growing the
result}: \emph{L-systems} (rewriting grammars generating branching structures),
\emph{cellular automata}, \emph{cellular encoding} (Gruau) and \emph{CPPN/HyperNEAT} for
neural networks (section~\ref{sec:neuro}). Their common contribution: a small genome, a large
phenotype, and a natural expression of symmetries and of repeated motifs.

\begin{defbox}[Novelty search and sparseness]
A radical idea (Lehman \& Stanley, 2011: \emph{Abandoning Objectives}): in deceptive problems
the objective function is misleading --- it leads into local optima. Therefore \textbf{let us
throw the goal away and reward novelty alone}.

A \emph{behaviour space} is defined (e.g.\ the final position of a robot in a maze) and the
novelty of an individual $x$ is measured by the sparseness of its neighbourhood,
\[
\rho(x) \;=\; \frac{1}{k}\sum_{i=1}^{k} \mathrm{dist}\big(x, \mu_i\big),
\]
where the $\mu_i$ are the $k$ nearest neighbours of $x$ in the current population
\textbf{and in a permanent archive}. Fitness $=\rho$; the original objective function
\emph{is not used at all}. On problems of the \uv{robot in a maze} type it tends to be more
successful than goal-directed optimisation. It is followed up by
\textbf{quality-diversity} (MAP-Elites): a grid of cells indexed by behaviour descriptors,
with the best solution found in each cell.
\end{defbox}

\paragraph{Interactive evolution.} The fitness function is replaced by a \textbf{human
being}: the user is presented with a population of candidates and picks the ones they like.
It is used where the quality criterion cannot be written down (aesthetics, design, music).
Examples: Dawkins's \emph{Biomorphs}, \textbf{Picbreeder} and \textbf{EndlessForms} (the
evolution of pictures and of 3D shapes by means of CPPNs). The limitation is slowness and
user fatigue --- small populations are necessary (9--20 individuals per screen) and few
generations. It was precisely the experience from Picbreeder (the most interesting objects
were never produced by deliberate search) that directly motivated novelty search.

\subsection{Exam summary}

\begin{itemize}
\item ES: real-valued vectors, mutation the main operator, random parent
  selection, deterministic environmental selection; an individual
  $= [\vec x, \vec\sigma]$.
\item $(\mu,\lambda)$ vs.\ $(\mu+\lambda)$: be able to list the differences
  (elitism, convergence, local optima, self-adaptation).
\item $(1+1)$-ES and the 1/5 rule: $p > 1/5 \Rightarrow$ increase $\sigma$,
  $p < 1/5 \Rightarrow$ decrease it; this is adaptive, not self-adaptive
  control.
\item Self-adaptation: $\sigma' = \sigma \exp(\tau'N(0,1) + \tau N_i(0,1))$, then
  $x' = x + \sigma' N(0,1)$; \textbf{first $\sigma$, then $x$}.
  The number of strategy parameters: 1 (isotropic), $n$ (uncorrelated),
  $n(n+1)/2$ (correlated, with rotation angles).
\item Recombination: local ($\rho=2$) / global ($\rho=\mu$),
  discrete (dominant) / intermediate (arithmetic).
\item CMA-ES: $\theta = (\vec m,\sigma,\mathbf C,\vec p_\sigma,\vec p_c)$,
  sampling from $\mathcal N(\vec m, \sigma^2\mathbf C)$, evolution paths,
  rank-1 and rank-$\mu$ updates, step control according to the length of
  $\vec p_\sigma$; invariant with respect to monotone transformations of fitness
  and to linear transformations of the space.
\item DE/rand/1/bin: the donor $\vec v = \vec x_a + F(\vec x_b - \vec x_c)$;
  binomial crossover with the threshold $C$ and the safeguard
  $j = I_{\mathrm{rand}}$; a duel of the offspring with its own parent
  (replacement only upon improvement).
\item Parameters: $F \in [0,2]$ (typically $0.5$--$0.9$), $C \in [0,1]$,
  $N \approx 10D$.
\item The geometric essence: the scale and the orientation of the steps are
  taken over from the current spread of the population $\Rightarrow$ automatic
  adaptation without strategy parameters.
\item Variants: rand/1, rand/2, best/1, best/2, current-to-best;
  \texttt{bin} vs.\ \texttt{exp} crossover.
\item Be able to demonstrate one step with concrete numbers (mutation $\to$
  crossover $\to$ selection).
\item Coevolution = contextual fitness + a dynamic landscape.
  \emph{Competitive}: predator--prey, host--parasite, Hillis's sorting networks.
  \emph{Cooperative}: decomposition into components, CCGA,
  modules $+$ blueprints.
\item Pathologies: cycling (intransitivity), disengagement, the Red Queen
  effect, mediocre stable states, overspecialisation, forgetting.
  The remedies: archives (hall of fame), shared fitness, measurement against a
  fixed set.
\item The prisoner's dilemma: $T>R>P>S$, $2R > T+S$; $D$ is dominant, $DD$ is a
  Nash equilibrium and the only outcome that is not Pareto-optimal; the iterated
  version $\Rightarrow$ TFT; the four properties of a successful strategy
  (niceness, provocability, forgiveness, non-enviousness) $+$ clarity as the
  fifth one; Pavlov in a noisy environment.
\item Diversity: fitness sharing, crowding, speciation, the island model,
  non-random mating.
\item Dynamic fitness: diploidy, hypermutation, random immigrants, maintaining
  diversity, the $(\mu,\lambda)$-ES.
\item Open-ended evolution: without a fixed goal; the hallmarks = sustained
  novelty, unbounded growth of complexity, new classes of entities, the
  evolution of evolvability. The systems Tierra, Avida, Polyworld, ECHO --- all
  of them get stuck sooner or later.
  Generative (indirect) representations: L-systems, cellular automata, cellular
  encoding, CPPNs.
\item Novelty search measures the sparseness $\rho(x)$ in the behaviour space
  with respect to the population \emph{and} a permanent archive, and does not
  use the objective function at all; it is followed up by quality-diversity and
  MAP-Elites (a grid of behaviour descriptors).
\item Interactive evolution: fitness = a human being (Biomorphs, Picbreeder,
  EndlessForms); the limitations are the small population and the user's
  fatigue.
\end{itemize}

%=====================================================================
\section{Swarm optimisation algorithms}
\label{sec:roje}
%=====================================================================

\begin{defbox}[Swarm intelligence]
An approach in which globally intelligent behaviour \textbf{emerges} from simple
local rules of many individuals without any central control. The features:
decentralisation, simpler individuals, indirect communication through the
environment (\emph{stigmergy}), robustness against the failure of an individual.
Unlike in EAs, there is usually no crossover, no mutation and no \uv{death} of
the individuals here --- the individuals move and remember.
\end{defbox}

\subsection{Particle swarm optimisation (PSO)}

Particle swarm optimisation (Eberhart and Kennedy, 1995) is inspired by flocks
of birds and schools of fish. An individual is a \emph{particle} --- a point in
$\mathbb{R}^n$ with a velocity. Neither crossover nor mutation is used; instead
the algorithm has a \textbf{memory}:
\begin{itemize}
\item $\vec p_i$ (\emph{pBest}) --- the best position particle $i$ has ever
  visited (the cognitive, individual memory),
\item $\vec g$ (\emph{gBest}) --- the best position found by the whole swarm
  (the social, collective memory).
\end{itemize}

\begin{defbox}[The PSO equations of motion]
\[
\vec v_i \;\gets\; \underbrace{w\,\vec v_i}_{\text{inertia}}
\;+\; \underbrace{c_1 r_1 (\vec p_i - \vec x_i)}_{\text{cognitive component}}
\;+\; \underbrace{c_2 r_2 (\vec g - \vec x_i)}_{\text{social component}},
\qquad
\vec x_i \;\gets\; \vec x_i + \vec v_i ,
\]
where $r_1, r_2 \sim U(0,1)$ (drawn separately for each component),
$c_1, c_2$ are the learning coefficients (in the original version
$c_1 = c_2 = 2$) and $w$ is the \emph{inertia weight}.
In the original 1995 version $w = 1$ (it is missing); Shi and Eberhart added it
and recommend a linear decrease, e.g.\ from $0.9$ to $0.4$: a large $w$ supports
exploration, a small one exploitation.
\end{defbox}

\begin{algorithm}[!ht]
\caption{PSO}
\begin{algorithmic}[1]
\State initialise the positions $\vec x_i$ and the velocities $\vec v_i$ at random; $\vec p_i \gets \vec x_i$
\Repeat
  \For{each particle $i$}
    \State evaluate $f(\vec x_i)$
    \If{$f(\vec x_i)$ is better than $f(\vec p_i)$} $\vec p_i \gets \vec x_i$ \EndIf
  \EndFor
  \State $\vec g \gets$ the best of $\{\vec p_i\}$
  \For{each particle $i$}
    \State update the velocity and the position according to the equations of motion
  \EndFor
\Until{the maximum number of iterations or the required precision is reached}
\end{algorithmic}
\end{algorithm}

\paragraph{A numerical example.} A particle in $\mathbb{R}^2$:
$\vec x = (1;\,2)$, $\vec v = (0.5;\,-0.5)$, $\vec p = (0;\,3)$,
$\vec g = (-1;\,1)$, $w = 0.7$, $c_1 = c_2 = 1.5$, $r_1 = 0.4$,
$r_2 = 0.8$.
\begin{align*}
\vec v' &= 0.7\,(0.5;\,-0.5) + 1.5\cdot 0.4\,\big((0;3)-(1;2)\big)
  + 1.5\cdot 0.8\,\big((-1;1)-(1;2)\big)\\
 &= (0.35;\,-0.35) + 0.6\,(-1;\,1) + 1.2\,(-2;\,-1)
  = (-2.65;\; -0.95),\\
\vec x' &= (1;\,2) + (-2.65;\,-0.95) = (-1.65;\; 1.05).
\end{align*}
The particle has thus overshot $\vec g$ --- typical PSO behaviour, and the
source of its exploration. That is why the velocity is bounded
($|v_j| \le v_{\max}$) or a \emph{constriction factor} (Clerc and Kennedy) is
used:
$\vec v \gets \chi[\vec v + c_1 r_1(\vec p - \vec x) + c_2 r_2 (\vec g - \vec x)]$,
where $\chi = 2/|2 - \varphi - \sqrt{\varphi^2 - 4\varphi}|$ for
$\varphi = c_1 + c_2 > 4$; for $c_1=c_2=2.05$ this gives
$\chi \approx 0.729$. This choice guarantees the convergence of the swarm.

\paragraph{Neighbourhood topologies.} Instead of a global $\vec g$ (the
\emph{gbest} topology, i.e.\ the complete graph) one can use \emph{lbest} ---
the best individual in a local neighbourhood:
\begin{itemize}
\item \textbf{A ring} (each particle has $k$ neighbours): information spreads
  slowly $\Rightarrow$ slower convergence, but better exploration and a smaller
  risk of getting stuck.
\item \textbf{A von Neumann grid}, \textbf{a star}, \textbf{random graphs}.
\item Empirically: gbest is faster on simple problems, lbest is more robust on
  multimodal ones.
\end{itemize}

\paragraph{Variants and the relation to EAs.}
Discrete/binary PSO (the velocity is interpreted as the probability of flipping
a bit via a sigmoid), PSO with several swarms, hybridisation with local search,
\emph{grammatical swarm}.

\begin{center}
\begin{tabular}{@{}p{5cm}p{5cm}p{5.2cm}@{}}
\toprule
 & \textbf{genetic algorithm} & \textbf{PSO} \\
\midrule
in common & random initialisation, fitness, stochastic updates &
  the same \\
individuals & are born and die & exist permanently, they move \\
memory & only in the population & explicit ($\vec p_i$, $\vec g$) \\
operators & crossover, mutation & none, only movement \\
information flow & two-way between the parents & one-way: from the better ones
  to the rest \\
\bottomrule
\end{tabular}
\end{center}

\subsection{Ant colony optimisation (ACO)}

\noindent\mimoq{of the swarm algorithms the EVA slides cover only PSO; the topic, however, speaks of swarm algorithms in the plural and ACO is their canonical second representative}

Ant colony optimisation (M.~Dorigo, 1991) imitates the way ants find a path: an
ant lays down a \emph{pheromone trail}, along a shorter path it returns sooner,
the pheromone accumulates there faster, and that attracts further ants ---
positive feedback. Moreover the pheromone \textbf{evaporates}, so the system
forgets obsolete information and does not get stuck.

\begin{defbox}[Stigmergy]
Indirect communication by means of changes in the environment. The ants do not
communicate with each other directly; the only channel is the pheromone on the
edges. A general principle of reactive multiagent architectures.
\end{defbox}

\paragraph{An illustration --- Steels's Mars explorer.} A swarm of robots is to
collect rock samples; the base emits a navigation signal (it suffices to follow
the gradient). A robot has \uv{crumbs} (radioactive markers) for indirect
communication. The rules, with priorities:
R1: if you see an obstacle, change direction; R2: if you are carrying a sample
and you are at the base, put it down; R3: if you are carrying a sample, go up
the gradient; R4: if you see a sample, pick it up; R5: otherwise move at random.
The second iteration adds stigmergy: R7: if you are carrying a sample and you
are not at the base, drop 2 crumbs and go up the gradient; R8: if you sense a
crumb, pick up 1 crumb and go down the gradient. What arises is an effective
collective collection from rich deposits, without the robots having either a map
or communication.

\paragraph{The general ACO scheme.} The problem is transformed into finding a
path in a weighted graph; the ants are agents constructing solutions along the
edges.
\begin{enumerate}
\item Every ant stochastically constructs a solution (a sequence of edges).
\item Optionally \emph{daemon actions} --- a centralised step, typically a local
  search improving the paths found.
\item The qualities of the solutions are compared.
\item The pheromones on the edges are updated (evaporation + deposition).
\end{enumerate}

\begin{defbox}[Ant System (AS) --- the transition rule]
Ant $k$ in city $i$ chooses the next city $j$ from the set of the not yet
visited ones $\mathcal{N}_i^k$ with probability
\[
p_{ij}^{k} \;=\;
\frac{\big[\tau_{ij}\big]^{\alpha}\,\big[\eta_{ij}\big]^{\beta}}
     {\sum_{l \in \mathcal{N}_i^k}\big[\tau_{il}\big]^{\alpha}\,\big[\eta_{il}\big]^{\beta}},
\qquad \eta_{ij} = \frac{1}{d_{ij}},
\]
where $\tau_{ij}$ is the amount of pheromone on the edge, $\eta_{ij}$ is the
heuristic \uv{visibility} (the reciprocal distance) and $\alpha, \beta$ are the
weights of the two influences (typically $\alpha = 1$, $\beta \in [2,5]$).
\end{defbox}

\begin{defbox}[Ant System --- the pheromone update]
\[
\tau_{ij} \;\gets\; (1-\rho)\,\tau_{ij} \;+\; \sum_{k=1}^{m} \Delta\tau_{ij}^{k},
\qquad
\Delta\tau_{ij}^{k} =
\begin{cases}
Q/L_k & \text{if ant } k \text{ used the edge } (i,j),\\
0 & \text{otherwise,}
\end{cases}
\]
where $\rho \in (0,1)$ is the evaporation rate and $L_k$ the length of the tour
of ant $k$. A shorter tour $\Rightarrow$ more pheromone.
\end{defbox}

\paragraph{A numerical example (the transition rule).} An ant stands in city $i$
and can go to the cities $A, B, C$ with distances $d = (5;\,10;\,2)$ and
pheromones $\tau = (2;\,4;\,1)$; $\alpha = 1$, $\beta = 2$.
The visibilities: $\eta = (0.2;\,0.1;\,0.5)$. The weights:
\[
w_A = 2 \cdot 0.2^2 = 0.08,\quad
w_B = 4 \cdot 0.1^2 = 0.04,\quad
w_C = 1 \cdot 0.5^2 = 0.25,\quad \textstyle\sum = 0.37 .
\]
The probabilities: $p_A = 0.216$, $p_B = 0.108$, $p_C = 0.676$.
Even though the edge to $B$ has the most pheromone, the short distance to $C$
(raised to the power $\beta = 2$) is what decides.

\paragraph{A numerical example (the update).} Let $\tau_{ij} = 2$,
$\rho = 0.5$, $Q = 1$, and let two ants with tour lengths $L_1 = 20$,
$L_2 = 25$ have used the edge:
\[
\tau_{ij} \gets 0.5 \cdot 2 + \left(\tfrac{1}{20} + \tfrac{1}{25}\right)
= 1 + 0.09 = 1.09 .
\]

\paragraph{ACS (Ant Colony System).} An improvement competitive with specialised
TSP solvers; inspired by Q-learning:
\begin{itemize}
\item \textbf{Global update by the best solution only} (elitism): pheromone is
  added only by the globally (or iteration-) best ant.
\item \textbf{Local update}: whenever an ant uses an edge, the pheromone on it
  is slightly decreased ($\tau \gets (1-\xi)\tau + \xi\tau_0$) --- the edge
  thereby becomes less attractive to the others and the variety of the tours
  grows.
\item \textbf{A pseudo-random selection rule}: with probability $q_0$ the best
  edge according to $\tau_{il}\eta_{il}^{\beta}$ is chosen deterministically
  (exploitation), otherwise the standard AS rule is used (exploration).
\end{itemize}

\paragraph{Other variants.}
\emph{Elitist AS} (the best tour contributes extra),
\emph{Rank-based AS} (the $N$ best contribute, with a weight according to the
rank),
\textbf{MAX--MIN AS} (the pheromone is kept in the interval
$[\tau_{\min}, \tau_{\max}]$, only the best one updates, initialisation to
$\tau_{\max}$ --- this prevents premature stagnation),
ACO for continuous optimisation (the space is discretised into subintervals and
the pheromone guides the search into better subspaces).

\paragraph{Properties and applications.} Decentralisation, emergence, indirect
communication; close to reactive agent architectures (Brooks's subsumption
architecture). A great advantage is the ability to \textbf{optimise in a dynamic
environment} --- when the graph changes, the pheromones adapt.
Applications: the TSP, vehicle routing, routing in networks, DNA sequencing,
protein folding, edge detection in images, scheduling.

\subsection{Other swarm algorithms}

\noindent\mimo{}

\begin{itemize}
\item \textbf{The artificial bee colony} (ABC; Karaboga 2005). The population of
  food sources = the solutions. Three roles:
  \emph{employed bees} (they search the neighbourhood of their source),
  \emph{onlookers} (they choose a source by the roulette wheel according to its
  quality and search its neighbourhood),
  \emph{scouts} (a source that has not improved for $limit$ attempts is
  abandoned and replaced by a random one --- a built-in restart).
\item \textbf{The firefly algorithm} (Yang, 2008). The individuals = fireflies;
  the brightness is proportional to fitness and decreases with distance as
  $I(r) = I_0 e^{-\gamma r^2}$. A less bright firefly moves towards a brighter
  one:
  $\vec x_i \gets \vec x_i + \beta_0 e^{-\gamma r_{ij}^2}(\vec x_j - \vec x_i)
  + \alpha \vec\varepsilon$. Thanks to the limited range of the attraction the
  swarm naturally splits into subgroups $\Rightarrow$ suitable for multimodal
  problems.
\item \textbf{The bat algorithm}, \textbf{cuckoo search}, \textbf{the grey wolf
  optimizer} and dozens of other \uv{metaphorical} algorithms. The criticism:
  many of them are merely renamed variants of PSO or ES and bring terminological
  confusion rather than new principles.
\end{itemize}

\subsection{Exam summary}

\begin{itemize}
\item PSO: the equations
  $\vec v \gets w\vec v + c_1r_1(\vec p - \vec x) + c_2r_2(\vec g - \vec x)$,
  $\vec x \gets \vec x + \vec v$; the inertia $w$ (decreasing 0.9 $\to$ 0.4),
  the cognitive and the social component, $c_1=c_2=2$; the bound $v_{\max}$ or
  the constriction $\chi \approx 0.729$; the gbest vs.\ lbest topology (ring,
  grid).
\item The difference between PSO and a GA: no operators, the particles have a
  memory ($\vec p_i$, $\vec g$), the information flows one way from the better
  ones.
\item ACO: the transition probability
  $p_{ij} \propto \tau_{ij}^{\alpha}\eta_{ij}^{\beta}$ with $\eta = 1/d$;
  the update $\tau \gets (1-\rho)\tau + \sum_k Q/L_k$ (evaporation + deposition
  proportional to the quality).
\item AS (the original one, everybody updates) vs.\ ACS (global update by the
  best one only, local update on traversal, the pseudo-random rule $q_0$)
  vs.\ MAX--MIN (bounds on the pheromone). Daemon actions = local search.
\item ABC (employed bees / onlookers / scouts with a limit = restart),
  firefly (the attraction decreases exponentially with the distance).
\item Keywords: stigmergy, emergence, decentralisation, positive feedback
  + evaporation, suitability for dynamic environments.
\end{itemize}


%=====================================================================
\section{Local search and memetic algorithms}
\label{sec:lokalni}
%=====================================================================

\subsection{Local search and hill climbing}

\begin{defbox}[Neighbourhood and local search]
For a solution space $X$ we define a \emph{neighbourhood}
$N(x) \subseteq X$ --- the set of solutions reachable from $x$ by one elementary
change (flipping a bit, swapping two cities, a 2-opt step). \emph{Local search}
iteratively moves to a neighbour according to a given rule. A \emph{local
optimum} with respect to $N$ is an $x$ such that no neighbour is better --- so
it depends on the definition of the neighbourhood!
\end{defbox}

\begin{defbox}[Hill climbing]
\textbf{1. Steepest ascent:} traverse the \emph{whole} neighbourhood and move to
the best neighbour if it is better than $x$; otherwise stop.\\
\textbf{2. First improvement:} traverse the neighbourhood and move to the first
better neighbour found (cheaper, often just as good).\\
\textbf{3. Stochastic hill climbing:} pick a random neighbour; accept it if it
is better (or with a probability proportional to the improvement).\\
\textbf{4. Random-restart hill climbing:} once stuck, start from a new random
position; remember the best solution found. As the number of restarts grows, the
probability of finding the global optimum converges to 1.
\end{defbox}

Weaknesses: getting stuck in a local optimum, on \emph{plateaux} (where all the
neighbours are the same --- \emph{sideways moves} help) and on ridges.
Advantages: simplicity, small memory, speed --- which is why it is ideal as an
improvement procedure inside an EA.

\subsection{Simulated annealing}

\begin{defbox}[Simulated annealing (SA)]
Kirkpatrick, Gelatt, Vecchi (1983); a physical analogy with the slow cooling of
a metal, during which the atoms take up an arrangement of minimal energy. The
algorithm is a hill climber that \textbf{also admits worsening steps} with a
probability that decreases with the size of the worsening and with the
temperature.
\end{defbox}

\begin{defbox}[The Metropolis criterion]
For minimisation, a candidate $y \in N(x)$ and $\Delta = f(y) - f(x)$:
\[
P(\text{accept } y) =
\begin{cases}
1, & \Delta \le 0 \quad (\text{an improvement}),\\[2pt]
e^{-\Delta / T}, & \Delta > 0 \quad (\text{a worsening}),
\end{cases}
\]
where $T > 0$ is the temperature.
\end{defbox}

\textbf{A numerical example.} A worsening of $\Delta = 2$. For $T = 1$ we get
$P = e^{-2} = 0.135$; for $T = 0.5$, $P = e^{-4} = 0.018$; for $T = 5$,
$P = e^{-0.4} = 0.67$. At the beginning (hot) the algorithm therefore behaves
almost like a random walk, at the end (cold) like a pure hill climber.

\begin{algorithm}[!ht]
\caption{Simulated annealing}
\begin{algorithmic}[1]
\State $x \gets$ an initial solution; $T \gets T_0$; $x^* \gets x$
\While{$T > T_{\min}$ and the termination condition is not met}
  \For{$k$ times (the length of the Markov chain at the given temperature)}
    \State pick $y \in N(x)$ at random; $\Delta \gets f(y) - f(x)$
    \If{$\Delta \le 0$ or $\mathrm{rand}() < e^{-\Delta/T}$} $x \gets y$ \EndIf
    \If{$f(x) < f(x^*)$} $x^* \gets x$ \EndIf
  \EndFor
  \State $T \gets \mathrm{cool}(T)$
\EndWhile
\Ensure $x^*$
\end{algorithmic}
\end{algorithm}

\paragraph{The cooling schedule.}
\begin{itemize}
\item \emph{Geometric:} $T_{k+1} = \alpha T_k$, $\alpha \in [0.8, 0.99]$ ---
  the most common one in practice.
\item \emph{Linear:} $T_k = T_0 - \beta k$.
\item \emph{Logarithmic:} $T_k = c/\log(k+1)$ --- the only one for which
  convergence has been proved.
\item Adaptive (reheating: the temperature is raised upon stagnation).
\end{itemize}

\begin{theorem}[Convergence of SA]
With a logarithmic schedule $T_k \ge c/\log(k+1)$ with a sufficiently large
constant $c$ (related to the height of the deepest \uv{barrier} in the
landscape), simulated annealing converges to the set of global optima with
probability 1.
\end{theorem}
In practice this schedule is unusably slow; SA is therefore theoretically
satisfactory, but faster heuristic schedules are used in practice. Formally SA
is an inhomogeneous Markov chain; for a fixed $T$ it converges to the Boltzmann
distribution $\pi(x) \propto e^{-f(x)/T}$, which for $T \to 0$ concentrates on
the global optima. The connection with EAs: the \emph{Boltzmann tournament}
(section~\ref{ssec:selekce}) carries the same idea over into selection.

\subsection{Tabu search}

\begin{defbox}[Tabu search (Glover, 1986)]
Local search with a \emph{short-term memory}: a \emph{tabu list} of recently
performed moves is kept (which is more efficient than storing whole solutions),
and those moves are forbidden for the duration of the \emph{tabu tenure}. In
each step one moves to the \textbf{best non-forbidden neighbour}, even if it is
worse than the current solution --- this is how the algorithm escapes from a
local optimum and avoids cycling.
\end{defbox}

\begin{itemize}
\item \textbf{Tabu tenure} $T$: static (a fixed number of iterations) or dynamic
  (drawn at random from a range).
\item \textbf{Aspiration criterion}: the tabu is broken if the move leads to a
  solution better than the best one found so far (otherwise the tabu could block
  progress).
\item \textbf{Medium-term memory}: rules steering the search into promising
  regions (intensification).
\item \textbf{Long-term / frequency memory}: a counter of how often each element
  of a solution has been used; frequently visited configurations are penalised
  (diversification), or one restarts from a rarely visited region.
\end{itemize}
Applications: the TSP (edge exchanges, \emph{ejection chains}), distribution
problems, DNA sequencing, scheduling. Pros: escape from local optima, no
cycling, it works in discrete as well as in continuous domains. Cons: many
parameters, higher computational cost (traversing the whole neighbourhood).

\subsection{Scatter search}

A population metaheuristic focused on \textbf{diversity}. It maintains a small
\emph{reference set} $R$ (on the order of tens of individuals) into which
individuals are selected according to a combination of quality and mutual
distance (the minimal distance of a newly added individual from the rest of $R$
is maximised).
The cycle: (1)~\emph{diversification} --- a clever initialisation,
(2)~\emph{subset generation} --- typically all the pairs from $R$,
(3)~\emph{combination} --- arithmetic (or permutation) crossover,
(4)~\emph{improvement} --- local search, mutation, tabu search,
(5)~\emph{updating $R$}. The reference set can be updated
\emph{incrementally} (one or a few individuals are added per generation), by a
\emph{complete renewal} every generation, or it can be built already in the
inner loop --- new individuals enter $R$ immediately and are recombined with
straight away. It is suitable for problems with a very expensive fitness
(e.g.\ the black-box optimiser OptQuest).

\subsection{Memetic algorithms}

\begin{defbox}[Meme and memetic algorithm]
A \emph{meme} (Dawkins, 1976) is a unit of cultural information spreading by
imitation --- a \uv{gene of culture}. A \emph{memetic algorithm} (Moscato, 1989)
is a hybrid method: \textbf{an evolutionary algorithm with a nested local
search} that improves the individual solutions (\uv{individual learning}).
Synonyms: hybrid EA, genetic local search, Lamarckian or Baldwinian EA, cultural
algorithms.
\end{defbox}

\begin{algorithm}[!ht]
\caption{Memetic algorithm}
\begin{algorithmic}[1]
\State initialise the population
\While{the termination condition is not met}
  \State evaluate the individuals in the population
  \State create the new population by the stochastic operators (selection, crossover, mutation)
  \State select a subset $O$ of individuals that will undergo improvement
  \For{each $x \in O$}
    \State perform individual learning of $x$ by means of a meme (a local searcher),
      with probability $p$ for a time $t$
    \State process the result in the \textbf{Lamarckian} or the \textbf{Baldwinian} way
  \EndFor
\EndWhile
\end{algorithmic}
\end{algorithm}

The essence: a local search nested inside the evolutionary cycle. The genetic
operators carry information \emph{between} the individual runs of the local
searcher, so the overall search is far more efficient than a repeated restart of
the local method (Krasnogor \& Smith). As a meme one can use hill climbing,
simulated annealing, tabu search, a combinatorial heuristic (2-opt) or a
gradient method (backpropagation when training a neural network).

\begin{defbox}[Lamarckism vs.\ Baldwinism]
What is to be done with an improved individual $x' = \mathrm{LS}(x)$?
\begin{itemize}
\item \textbf{The Lamarckian way:} $x'$ is inserted into the population together
  with its fitness --- \emph{the genotype is changed according to the
  phenotype}. From a Darwinian point of view this is \uv{not kosher}, but in
  practice it is faster.
\item \textbf{The Baldwinian way:} the individual $x$ is assigned the
  \emph{fitness} of the improved $x'$, but the genotype of $x$ remains
  \emph{unchanged}. This is Darwinistically correct; the fitness is interpreted
  as the \uv{potential} of the individual (the \emph{Baldwin effect}: the
  ability to learn favours an individual and evolution subsequently
  \uv{catches up} by itself --- genetic assimilation).
\end{itemize}
Lamarckism converges faster, but it reduces diversity more and may distort the
landscape; Baldwinism preserves diversity and smooths the fitness landscape, but
it is slower.
\end{defbox}

\paragraph{Design questions.} Which individuals should be improved (all / the
elite only / a random fraction $p$)? For how long ($t$ steps)? How strongly (the
depth of the local search)? Too much local search means unnecessary cost and a
loss of diversity; too little brings no benefit.

\paragraph{Memetic computing.} A generalisation: the memes are not known in
advance but are themselves the object of the search.
\begin{itemize}
\item \textbf{Multi-meme MA:} the meme is encoded \emph{as part of the genotype}
  of an individual; it is inherited, it is crossed (the offspring inherits the
  meme of the better parent, or at random) and it is used to improve that
  individual.
\item \textbf{Meta-Lamarckian MA:} the memes have \emph{their own population}
  and coevolve with the individuals; a meme is assigned to an individual for the
  improvement, the success of the individual is the reward for the meme, and a
  second evolutionary algorithm runs over the memes.
\item The \emph{memetic space} as a counterpart of the phenotype space: memes
  can be chosen by heuristics, made to compete according to their success, split
  into subpopulations, or controlled by a meta-memetic algorithm.
\end{itemize}

\subsection{Parameter tuning and control}

This is closely related: an EA has many parameters (the population size, $p_C$,
$p_M$, the tournament size, the degree of elitism). They are not independent and
an exhaustive search is not possible.
\begin{itemize}
\item \textbf{Tuning} --- the parameters are determined \emph{before} the run
  (trial and error, grid search, a meta-EA, iterated racing --- irace).
\item \textbf{Control} --- the parameters change \emph{during} the run:
  \begin{itemize}
  \item A \emph{fixed strategy} (deterministic): a dependence on time only
    (e.g.\ $\sigma(t) = 1 - 0.9\,t/T$, a decrease of $p_M$ by 0.1~\% per
    generation); the stochastic variant = simulated annealing.
  \item \emph{Adaptive}: according to feedback from the run (the quality of the
    solutions, the variance of fitness, the diversity) --- e.g.\ the 1/5 rule,
    hypermutation upon a drop in fitness, adaptive refinement of the
    representation (increasing the number of bits per real number if the
    diversity, measured by the Hamming distance of the best and the worst
    individual, drops).
  \item \emph{Self-adaptive}: the parameters are part of the genome and evolve
    (ES, EP, multi-meme MA).
  \end{itemize}
\item The \emph{population size} can also be controlled indirectly: every new
  individual is assigned a \uv{lifetime} proportional to its fitness, and it
  dies once that expires. Better individuals live longer, and the population
  size is a derived dynamic quantity.
\end{itemize}

\subsection{Exam summary}

\begin{itemize}
\item Hill climbing: steepest ascent / first improvement / stochastic / with
  random restarts; the problems: local optima, plateaux, ridges.
\item SA: the Metropolis criterion $P = \min(1, e^{-\Delta/T})$; the cooling
  schedule (geometric in practice, logarithmic for the convergence proof); for a
  fixed $T$ it converges to the Boltzmann distribution $\propto e^{-f/T}$.
\item Tabu search: a tabu list of moves, the tenure, the aspiration criterion,
  medium-term and long-term (frequency) memory; one always moves to the best
  non-forbidden neighbour, even if it is worse.
\item Scatter search: a reference set combining quality and diversity,
  arithmetic combinations + improvement.
\item Memetic algorithm = EA + local search over the individuals;
  \textbf{Lamarckism} (it changes the genotype, faster) vs.\
  \textbf{Baldwinism} (it changes only the fitness, it preserves diversity).
  Memetic computing: memes in the genome or in a coevolving population.
\item Parameter tuning vs.\ control; fixed / adaptive / self-adaptive
  strategies.
\end{itemize}


%=====================================================================
\section{Applications of evolutionary algorithms}
\label{sec:aplikace}
%=====================================================================

The last sentence of the topic enumerates four application areas; each forms one section of
this chapter. What they have in common is that in each of them what decides is the
\textbf{choice of representation and operators}, not the choice of a \uv{better}
evolutionary algorithm.

\subsection{Evolution of rule-based and expert systems}
\label{sec:lcs}

\subsubsection{Michigan vs.\ Pittsburgh}

The goal is to learn a \textbf{set of rules} of the form
\texttt{IF condition THEN action} from training data or from interaction with an
environment (knowledge discovery, expert systems, robot control, reinforcement
learning). There are two basic approaches:

\begin{center}
\begin{tabular}{@{}p{2.8cm}p{6.2cm}p{6.2cm}@{}}
\toprule
 & \textbf{Michigan} (Holland) & \textbf{Pittsburgh} (Smith) \\
\midrule
individual & \textbf{a single rule} & \textbf{a whole set of rules} \\
solution & the whole population together & a single individual \\
fitness & the strength/accuracy of the rule (the reward has to be distributed) &
  the quality of the whole system (straightforward) \\
operators & standard, evolution runs only occasionally and on a part of the
  population &
  complex, dozens of operators over rule sets and over individual rules \\
advantages & on-line learning, a small space, natural for RL & an unambiguous
  fitness, no credit assignment \\
disadvantages & credit assignment; the rules have to cooperate, yet they compete
  at the same time &
  large individuals, expensive evaluation, a risk of bloat \\
\bottomrule
\end{tabular}
\end{center}

\subsubsection{Learning classifier systems (Michigan)}

\begin{defbox}[LCS (learning classifier system)]
A system in which a population of simple rules (\emph{classifiers}) forms an
expert or control system as a whole. A classifier has a left-hand side (the
condition) in the form of a string over $\{0,1,*\}$ that is matched against the
input, a right-hand side (an action or a classification class) and a
\textbf{weight / strength} reflecting its success, which serves as its fitness.
Evolution does not proceed generationally, but continuously and only on a part
of the population.
\end{defbox}

\paragraph{The architecture of an LCS.} The key point to realise is that
\emph{the expert system is the whole population}, not an individual. The system
forms a loop:
\begin{enumerate}
\item The \textbf{detectors} convert the state of the environment into an input
  message.
\item The message is stored in the \textbf{message buffer} and matched against
  the left-hand sides of all the rules $\Rightarrow$ the \emph{match set}.
\item Among the matching rules an \textbf{auction/competition} for the right to
  act takes place (according to the strength of the rule).
\item The winning rule writes its right-hand side either into the buffer (an
  internal message) or, through the \textbf{effectors}, into the environment as
  an action.
\item The environment returns a \textbf{reward}; this is distributed by the
  \emph{bucket brigade} mechanism among the rules that led to it.
\item Occasionally, and only on a part of the population, a \textbf{GA} runs and
  renews the rules.
\end{enumerate}

\paragraph{The problem of reactivity.} A purely reactive system has no internal
memory. Holland therefore introduced \emph{messages}: besides an action, the
right-hand side of a rule may write an internal message into the buffer, and the
left-hand side has \emph{receptors} for catching it. Chaining of rules arises in
this way --- the system acquires an internal state (and thereby computational
power), but at the same time it brings the problem of how to distribute the
reward over the whole chain.

\begin{defbox}[Bucket brigade]
The mechanism for distributing the reward in an LCS. A rule that wants to be
applied must \uv{pay} a part of its strength to the rule that activated it (as
if it were buying the right to act). The reward from the environment goes to the
last link of the chain and gradually \uv{flows} back to the rules that led to
it. In practice it is difficult to balance this economy of rules, which is why
it is hardly ever used today.
\end{defbox}

\paragraph{ZCS (Wilson, 1994) --- the \uv{zero-level} LCS.} A simplification: no
internal messages, no complex redistribution.
\begin{itemize}
\item $P$ --- the population of rules; the input $x$ from the detectors.
\item $M$ --- the \emph{match set}: the rules whose condition matches $x$.
\item The action $a$ is selected from $M$ by the roulette wheel according to the
  strengths of the rules; $A \subseteq M$ --- the \emph{action set}: the rules
  advocating the chosen action.
\item \textbf{Reinforcement:} a fraction $\beta$ of the strength of the rules in
  $A$ is taken away into a \uv{bucket} $B$; after a reward $r$ is received, every
  rule in $A$ has $\beta r/|A|$ added to it and the rules of the previous action
  set $A_{-1}$ have $\gamma B/|A_{-1}|$ added to them. The rules from
  $M \setminus A$ (they matched, but advocated a different action) are taxed.
\item \textbf{The cover operator:} if no rule exists for the current input, one
  is generated ad hoc (stars are added at random into the condition and the
  action is chosen at random).
\end{itemize}

\paragraph{XCS (Wilson, 1995) --- accuracy-based fitness.}
The weaknesses of ZCS: it does not evolve a \emph{complete} system of rules
covering all situations; rules at the beginning of chains are rewarded rarely
and die out; rules leading to small but correctly predicted rewards die out as
well. The cause is that the strength served both as the fitness for the GA and
as the criterion for choosing the action. Wilson's idea: the \emph{prediction}
should estimate the reward, but \emph{evolution} should prefer
\textbf{reliable (accurate)} classifiers.

\begin{defbox}[XCS --- a classifier as a quintuple]
$(c, a, p, \epsilon, F)$: the condition, the action, the reward prediction $p$,
the prediction error $\epsilon$, the fitness $F$. The accuracy is computed as a
decreasing function of the error,
$\kappa = \alpha\,(\epsilon/\epsilon_0)^{-\nu}$ (for $\epsilon > \epsilon_0$,
otherwise $\kappa = 1$; $\alpha, \epsilon_0, \nu$ are parameters);
the fitness is the \emph{relative} accuracy within the action set.
The prediction and the error are updated by the rules of Q-learning.
\end{defbox}

A run of XCS: according to the input the match set $M$ is assembled and split
into action sets $A_i$ according to the actions; for each action the payoff is
predicted as the fitness-weighted average of the predictions; the action is
chosen either by the maximum (exploitation) or by the roulette wheel
(exploration); it is performed; the reward is distributed into the action set of
the previous step (an update of $p$ using $r + \gamma \max_a p_a$, then
$\epsilon$, then $F$); the GA runs as a \textbf{niche} GA --- only inside the
action set, not over the whole population.

The result: XCS links LCS with reinforcement learning (Q-learning as the credit
assignment mechanism), evolves a \textbf{complete and minimal} input--output map
(not only an optimal policy, but also a compact, comprehensible description) and
has a number of practical applications (knowledge discovery, control,
bioinformatics).

\paragraph{Today's family of LCS.} Depending on the task one distinguishes:
\emph{XCS} (reinforcement learning, reward prediction),
\textbf{UCS} (\emph{sUpervised Classifier System}) --- the variant for
supervised learning, where the correct class is known, so that classification
accuracy is measured directly instead of a reward prediction,
and \emph{XCSF} for the approximation of continuous functions (the right-hand
side is a linear model, not a constant). Nowadays the conditions are commonly
written not only over $\{0,1,*\}$ but also with intervals for real-valued
attributes, which makes LCS a fully-fledged knowledge discovery tool --- with
the advantage that the result is a \textbf{readable set of rules}.

\subsubsection{The Pittsburgh approach and the GIL system}

An individual = a whole set of rules, i.e.\ a complete classification system.
The evaluation is more complicated (rule priorities, conflicts, false positive
and false negative answers) and there is a whole range of operators, at several
levels. The emphasis is on a rich domain representation (sets, enumerations,
intervals).

\textbf{GIL} is the classical representative: it is meant for binary
classification, and an individual classifies implicitly into one class (the
rules have no right-hand side). An individual is a \emph{disjunction of
complexes}, a complex is a \emph{conjunction of selectors} (one for each
variable) and a selector is a \emph{disjunction of values} from the domain of
the given variable, encoded by a bit mask. For example
\[
\big((X{=}A_1) \wedge (Z{=}C_3)\big) \vee \big((X{=}A_2) \wedge (Y{=}B_2)\big)
\;\equiv\;
[\,001|11|0011 \;\vee\; 010|10|1111\,].
\]
The operators work at three levels:
\begin{itemize}
\item the \emph{individual}: exchange of rules, copying a rule, generalising a
  rule, deleting a rule, specialising a rule, including a positive example;
\item the \emph{complex}: splitting a complex on one selector, generalising a
  selector (replacing it by $11\ldots1$), specialising it, including a negative
  example;
\item the \emph{selector}: mutation $0 \leftrightarrow 1$, widening $0 \to 1$,
  narrowing $1 \to 0$.
\end{itemize}
These operators correspond directly to the operations of inductive learning
(generalisation and specialisation), so GIL is in fact an evolutionarily driven
inductive algorithm.

\subsubsection{A reminder: classification metrics}

From the confusion matrix (TP, FP, FN, TN) one computes
$\text{accuracy} = (TP+TN)/(TP+TN+FP+FN)$,
$\text{sensitivity} = TP/(TP+FN)$ and $\text{specificity} = TN/(TN+FP)$.
The fitness of a rule-based system is usually a combination of accuracy, coverage and
simplicity (the number of rules) --- so it is naturally a multi-objective problem.

\subsection{Neuroevolution}
\label{sec:neuro}

\begin{defbox}[Neuroevolution]
\uv{Neuroevolution is a method for modifying neural network weights, topologies,
or ensembles in order to learn a specific task. Evolutionary computation is used
to search for network parameters that maximize a fitness function that measures
performance in the task. Compared to other neural network learning methods,
neuroevolution is highly general --- it allows learning without explicit
targets, with non-differentiable activation functions, and with recurrent
networks.} (Miikkulainen, 2017)
\end{defbox}

What can be evolved: the weights, further parameters (activation functions,
learning constants), the architecture (the topology, the connections), the
architecture and the weights at the same time, or the parameters of the learning
algorithm itself. The main domain is \textbf{reinforcement learning} and
robotics, where target outputs for supervised learning are not available.

\subsubsection{Evolution of the weights}

Straightforward: the weights are encoded into a vector of fixed length and a
real-valued GA, an ES or DE with the standard operators is used. Properties:
\begin{itemize}
\item Slower than specialised gradient methods (backpropagation) for supervised
  learning.
\item \textbf{+} Easy parallelisation, no gradient needed (it works for
  non-differentiable and recurrent networks), usable for RL.
\item \textbf{The problem of competing conventions} (the permutation problem):
  the same network function corresponds to many different genotypes
  (permutations of the hidden neurons), so crossing two functionally identical
  parents often produces a non-functional offspring.
\item A recent revival of interest: the evolution of weights is competitive even
  for large networks if evaluation on minibatches is used.
\end{itemize}

\subsubsection{Evolution of the architecture}

The fitness of an architecture = an estimate of the performance of the network:
the network has to be built, initialised and trained (by backpropagation or by
another EA), preferably several times --- which is why the fitness evaluation is
expensive and noisy.

\begin{itemize}
\item \textbf{Direct encoding}: the connection structure as a binary matrix; an
  individual is either the linearised matrix (standard operators) or the matrix
  itself (special 2D operators). Simple, but it does not scale (the size of the
  genome grows quadratically).
\item \textbf{Grammar encoding} (Kitano, 1990): the connection matrix is
  generated by a simple 2D formal grammar and the \emph{rules of the grammar}
  are evolved. Compact (logarithmic), it allows repeated motifs, but it is
  indirect and harder to tune.
\item \textbf{Cellular encoding} (Gruau): an individual is a \textbf{GP program}
  describing how to \emph{grow} the network from a single cell by means of the
  operations: add a neuron, split a neuron serially, split it in parallel,
  reconnect a synapse, change a weight. A developmental approach.
\item \textbf{Simulated growth} of connections in a 2D plane --- early attempts
  in evolutionary robotics; it did not scale.
\end{itemize}

\paragraph{GNARL (Angeline et al., 1993).} Evolutionary programming over graphs:
the input and output neurons are fixed, the hidden neurons and the connections
are evolved; the architecture and the weights \emph{at the same time}. Two kinds
of mutation: \emph{structural} (add a neuron without connections, add a
connection with zero weight, delete a neuron, delete a connection) and
\emph{parametric} (biased mutation of a weight). The strength of the mutation is
controlled by a \emph{temperature} (a simulated annealing approach); structural
changes should be small and not too destructive. Selection: the better half of
the population becomes the parents. Crossover is not used (EP).

\subsubsection{NEAT}

\begin{defbox}[NEAT (NeuroEvolution of Augmenting Topologies,
K.~Stanley, 2002)]
The network is represented as a \textbf{list of edges}; every edge gene
contains: the input and the output neuron, the weight, an \emph{enabled/disabled}
flag and a \textbf{global innovation number}, which is assigned when the given
structural novelty first arises.
\end{defbox}

Three key ideas:
\begin{enumerate}
\item \textbf{Innovation numbers solve the problem of competing conventions.}
  Only genes with \emph{matching innovation numbers} are crossed (so they have a
  common evolutionary origin); genes that are surplus in one of the parents
  (\emph{disjoint} and \emph{excess}) are taken over from the fitter parent.
  Crossover thus does not disrupt the structure --- no \uv{headless chicken}.
\item \textbf{Speciation (niching) protects innovations.} On the vectors of
  innovation numbers a similarity measure is defined:
  \[
  \delta \;=\; \frac{c_1 E}{N} + \frac{c_2 D}{N} + c_3 \overline{\Delta W},
  \]
  where $E$ (\emph{excess}) is the number of genes whose innovation number lies
  \emph{beyond} the highest innovation number of the other genome, $D$
  (\emph{disjoint}) is the number of genes missing in the other parent
  \emph{inside} the common range of innovation numbers, $\overline{\Delta W}$ is
  the average difference of the weights of genes with \emph{matching} innovation
  numbers and $N$ is the size of the larger genome (for normalisation). Networks
  with $\delta$ below a threshold form one \emph{species};
  fitness is computed as shared (relative within the species). A new topology
  thus has time to \uv{fine-tune its weights\,} before established solutions
  crowd it out --- otherwise a structural novelty that initially worsens the
  fitness would die out immediately.
\item \textbf{Gradual growth of complexity} (\emph{complexification}): the
  population starts from minimal networks (inputs and outputs only) and grows
  only through the mutations \emph{add a connection} and \emph{add a neuron}
  (which splits an existing connection: the original one is disabled and two new
  ones arise). The search therefore starts in the space of the lowest dimension
  and increases it only when it pays off.
\end{enumerate}

\subsubsection{HyperNEAT and CPPNs}

\noindent\mimoq{the slides give it one line, CPPNs not at all}

\textbf{HyperNEAT} extends NEAT by two ideas. \emph{The substrate}: the weights of the
network are not encoded as a vector of numbers, but \emph{are generated by a function}
$S : \mathbb{R}^4 \to \mathbb{R}$ which receives the coordinates of two neurons in the plane
and returns the weight of their connection (hence \uv{hyper}). \emph{A CPPN}
(compositional pattern-producing network): the function $S$ is represented by a network in
the form of a DAG whose nodes compute various functions from a dictionary (sigmoid, Gaussian,
sine, absolute value); this CPPN is evolved by NEAT itself. The benefit: representing the
weights by a function makes it possible to express \textbf{symmetries and repeated motifs}
and to generate networks with millions of connections from a small genome.

\subsubsection{Neuroevolution in the era of deep networks}

\noindent\mimo{}

\textbf{NAS} (neural architecture search) looks for the architecture of deep networks and is
described along three axes: the \emph{search space} (a sequence of layers, a DAG of repeated
cells, or a subarchitecture of one large \uv{supernetwork}), the \emph{search strategy}
(evolutionary algorithms, random search as a surprisingly strong baseline, Bayesian
optimisation, reinforcement learning, gradient-based DARTS) and \emph{performance estimation}
(shortened training, weight sharing). Concrete systems: \textbf{ENAS} (weight sharing among
the candidates), \textbf{DeepNEAT} (a node of the chromosome is a whole layer) and
\textbf{CoDeepNEAT} (coevolution of a population of modules and a population of
\uv{blueprints}). A separate line is \textbf{ES for RL} (OpenAI 2017): a simple evolution
strategy \emph{with mutations only} over millions of weights, trivially parallelisable ---
the nodes exchange only the seeds of the random number generator.

\subsection{Combinatorial problems and constraint handling}
\label{sec:kombi}

Evolutionary algorithms are a popular tool for NP-hard combinatorial problems:
they do not guarantee the optimum, but they provide a good solution in a
reasonable time and adapt easily to the additional (possibly quite ugly)
constraints of a real problem.

\subsubsection{The knapsack problem}

\textbf{The task.} A knapsack of capacity $C_{\max}$, $N$ items with values
$v_i$ and volumes $c_i$. Maximise $\sum_i x_i v_i$ subject to
$\sum_i x_i c_i \le C_{\max}$, $x_i \in \{0,1\}$.

\textbf{The encoding} is trivial --- a bit map, e.g.\ $0110010$ means
\uv{take items 2, 3 and 6}. The standard operators (crossover, bit-flip
mutation) work. \textbf{The problem is the fitness}: individuals may violate the
capacity. We have two conflicting criteria: to maximise $\sum v_i$ and at the
same time to satisfy $\sum c_i \le C_{\max}$.

\subsubsection{Three general strategies for constraints}

\begin{defbox}[Constraint handling in EAs]
\textbf{1. Penalty}: infeasible solutions stay in the population, but their
fitness is reduced, e.g.
$f'(x) = \sum_i x_i v_i - \alpha \max\big(0, \sum_i x_i c_i - C_{\max}\big)$.
The choice of $\alpha$ is crucial: too small $\Rightarrow$ the population fills
up with infeasible solutions, too large $\Rightarrow$ the possibility of passing
through the infeasible region towards a promising solution is lost. Dynamic
($\alpha$ growing in time) or adaptive penalties are also used.\\
\textbf{2. Repair}: an infeasible solution is repaired (for the knapsack: throw
out the items with the worst ratio $v_i/c_i$ until it fits). The repair can be
performed only for the computation of fitness (the Baldwinian way) or written
permanently into the genotype (the Lamarckian way).\\
\textbf{3. Decoder} / closed operators: the encoding is designed so that
\emph{every} genotype represents a feasible solution. For the knapsack: a bit
1 means \uv{add the item if it still fits}, otherwise skip it. Elegant and
effective; the drawback is the ambiguity of the mapping (several genotypes give
the same phenotype) and its poorer locality.
\end{defbox}

A fourth possibility: to regard the violation of the constraint as a
\textbf{second criterion} and to use a multi-objective EA
(section~\ref{sec:moea}) --- the result is then a whole Pareto front of
trade-offs between the quality of the solution and the degree of constraint
violation.

\subsubsection{The travelling salesman problem}

The TSP: $N$ cities, find the shortest round tour. The fitness is trivial (the
length of the tour); \textbf{the problem is the encoding and above all the
crossover}.

\paragraph{Representations.}
\begin{itemize}
\item \textbf{Adjacency}: at position $i$ there is city $j$ precisely if there
  is an edge $i \to j$. For instance, $(248397156)$ corresponds to the tour
  $1{-}2{-}4{-}3{-}8{-}5{-}9{-}6{-}7$. Every tour has a unique representation,
  but some lists are not valid tours. Classical crossover does not work; what is
  interesting is that schemata make sense here (e.g.\ $(*3*\ldots)$ = all the
  tours with the edge 2--3). \emph{Do not use.}
\item \textbf{Ordinal} (\emph{buffer}): a list of cities is maintained, gene $i$
  is the \emph{position} of the city in the current list, and after being used
  the city is removed from the list. For the list $(123456789)$ the tour
  $1{-}2{-}4{-}3{-}8{-}5{-}9{-}6{-}7$ is encoded as $(112141311)$. The
  motivation was to make standard one-point crossover possible --- and it does
  produce valid tours, but the result has nothing in common with the parents.
  \emph{Also not to be used.}
\item \textbf{Path (permutation)}: the tour
  $5{-}1{-}7{-}8{-}9{-}4{-}6{-}2{-}3$ is $(517894623)$. The most natural one; it
  requires special operators --- PMX, OX, CX, ER
  (section~\ref{sec:reprez}).
\item \textbf{Matrix}: a binary matrix in which a $1$ at position $(i,j)$ means
  either the edge $i \to j$ or (more often) that $i$ precedes $j$. Special
  operators: \emph{conjunction} (a bitwise AND of the two parents and a random
  completion of the missing edges), \emph{disjunction} (a split into quadrants,
  two of which are discarded, the contradictions are removed and the rest is
  filled in at random).
\end{itemize}

\paragraph{Initialisation.} Random permutations, or constructive heuristics:
\emph{nearest neighbour} (start in a random city and always go to the nearest
unvisited one) and \emph{edge insertion} (for a partial tour $T$ choose the
nearest city $c$ outside $T$, find the edge $k$--$j$ in $T$ minimising
$d(k,c) + d(c,j) - d(k,j)$, and replace that edge by the pair $k$--$c$,
$c$--$j$).

\paragraph{Mutation.} Inversion of a segment (the most important one --- it
corresponds to a 2-opt step), inserting a city elsewhere, shifting a subtour,
swapping two cities, swapping subtours, the heuristic mutations 2-opt and 3-opt
(take two edges and reconnect four cities differently). A combination of an ES
or a GA with \uv{clever mutations} of the 2-opt kind is very effective in
practice --- de facto a memetic algorithm.

\subsubsection{SAT}

\noindent\mimoq{the slides do the knapsack, the TSP and scheduling; SAT comes from Michalewicz}

\textbf{The task.} A formula $f: \mathbb{B}^n \to \mathbb{B}$ in CNF; find an assignment $x$
with $f(x)=1$. 2-SAT is in P, 3-SAT and above is NP-complete.

\textbf{Fitness.} $f$ itself is unusable (it takes only the values 0/1, so the landscape is a
\uv{needle in a haystack}). The standard is the \textbf{number of satisfied clauses}; the
improvements are weighting the problematic clauses and a \emph{refining function} --- a
regularisation term that distinguishes individuals with the same number of satisfied clauses
and thus helps off the plateaux.

\textbf{The representation} is usually a plain bit string. The \emph{real-valued} variant is
interesting: the variable $x_j$ is replaced by a real $y_j$ (true $=1$, false $=-1$), the
literals are rewritten as $x \mapsto (1-y)^2$ and $\neg x \mapsto (1+y)^2$, and the expression
is \emph{minimised}. The value of a literal is therefore a \emph{penalty}, and that is why a
\textbf{disjunction} (a clause) is rewritten as a \textbf{product} and a \textbf{conjunction}
(the formula) as a \textbf{sum}. The classical local heuristic is \textbf{WalkSAT}; the
combination EA + WSAT is a typical memetic solution.

\subsubsection{Scheduling and production planning}

\textbf{A school timetable.} An individual is a timetable in the form of a matrix (rows =
teachers, columns = classes, values = subject codes). Mutation shuffles the subjects,
crossover exchanges the better rows between individuals. The fitness combines the quality of
the individual rows and soft criteria (gaps, uniformity). \textbf{Hard constraints} (a
teacher cannot teach two classes at once) have to be respected \emph{directly in the
operators} --- otherwise far too many infeasible individuals arise.

\textbf{Job shop scheduling.} Products consist of parts, for each part there are several
production plans on the machines, and resetting a machine costs time; the fitness is the
total production time (the \emph{makespan}). The encoding is critical: a
\emph{permutation-based} individual (only the order of the products, the rest determined by a
decoder) allows the TSP crossovers to be used, but turns out to be not very effective --- the
difficult part is solved by the decoder. A \emph{direct} encoding of the whole plan is more
powerful but requires specialised operators.

\subsection{Multi-objective optimisation}
\label{sec:moea}

\subsubsection{Dominance and the Pareto front}

Instead of a single fitness we have a vector of criteria $f_1,\dots,f_k$ (for
simplicity we minimise everything). The criteria are usually in conflict (cost
vs.\ quality, accuracy vs.\ the size of the model), so there is no single best
solution.

\begin{defbox}[Dominance]
For individuals $x, y$ we define:
\begin{itemize}
\item $x$ \textbf{weakly dominates} $y$ ($x \preceq y$) if and only if
  $f_i(x) \le f_i(y)$ for all $i = 1,\dots,k$.
\item $x$ \textbf{dominates} $y$ ($x \prec y$) if and only if $x \preceq y$ and
  there exists a $j$ with $f_j(x) < f_j(y)$.
\item $x$ and $y$ are \textbf{incomparable} if neither $x \prec y$ nor
  $y \prec x$.
\end{itemize}
The relation $\prec$ is a \emph{strict} partial order --- it is irreflexive and
transitive, but not total: that is why in general there is a whole set of
mutually incomparable optima.
\end{defbox}

\begin{defbox}[Pareto optimality and the Pareto front]
A solution $x^*$ is \emph{Pareto-optimal} if it is not dominated by any feasible
solution. The set of all Pareto-optimal solutions in the space of the variables
is the \emph{Pareto set}, and its image in the space of the criteria is the
\textbf{Pareto front}. An \emph{approximation} of the Pareto front by a
population is the goal of multi-objective EAs.
\end{defbox}

The algorithm has a twofold goal: (a) \textbf{convergence} --- to get as close
as possible to the true front; (b) \textbf{diversity} --- to cover the front
evenly over its whole extent. This is exactly why EAs are ideal for this task:
they work with a population and can therefore return a whole approximation of
the front in a single run.

\textbf{Example.} The points $A(1,10)$, $B(2,7)$, $C(4,5)$, $D(8,2)$, $E(5,8)$
(both criteria minimised). $E$ is dominated by $C$ ($4 \le 5$, $5 \le 8$, at
least one of them strict) --- so $E$ does not belong to the front. The points
$A, B, C, D$ are mutually incomparable and form the Pareto front.

\subsubsection{Scalarisation --- the simple solutions}

\begin{itemize}
\item \textbf{Linear scalarisation}: $f = \sum_i w_i f_i$, followed by a
  standard single-objective EA (referred to as an SOEA in the MOEA context). The
  drawbacks: we do not know how to choose the weights; one run gives one point
  of the front; and, crucially, \textbf{non-convex parts of the front can never
  be found by linear scalarisation}.
\item \textbf{$\varepsilon$-constraint}: minimise $f_1$ subject to the
  constraints $f_i \le \varepsilon_i$ for $i \ge 2$. It can find non-convex
  parts as well, but the constants $\varepsilon_i$ have to be chosen and a
  constrained problem has to be solved.
\item \textbf{Chebyshev scalarisation}: minimise
  $\max_i \lambda_i |f_i(x) - z_i^*|$, where $z^*$ is a reference (ideal) point.
  Unlike the linear one, it can reach any point of the front, including the
  non-convex parts.
\end{itemize}

\subsubsection{The historical algorithms}

\paragraph{VEGA (Schaffer, 1985)} --- one of the first MOEAs. The population of
$N$ individuals is sorted for each criterion $f_i$ and the $N/k$ best according
to $f_i$ are selected; these subsets are merged, crossed and mutated.
The drawbacks: it is difficult to maintain diversity and it converges to the
extreme points optimal for the individual criteria (the \uv{middles} of the
front are lost).

\paragraph{NSGA (1994)} --- the first use of dominance directly as the fitness.
The population is split into fronts: $F_1$ = the non-dominated individuals of
$P$, $F_2$ = the non-dominated individuals of $P \setminus F_1$,
$F_3$ = the non-dominated ones of $P \setminus (F_1 \cup F_2)$, and so on. The
individuals of the first front receive a \uv{dummy} fitness, which is divided by
a \emph{niche factor} $\sum_j sh(i,j)$, where
$sh(i,j) = 1 - (d(i,j)/d_{\mathrm{share}})^2$ for $d(i,j) < d_{\mathrm{share}}$
and 0 otherwise. The second front receives a fitness lower than the worst
individual of the first front, again divided by the niche factor, and so on.
The drawbacks: the need to set $d_{\mathrm{share}}$, the absence of elitism, and
the computational complexity $O(kN^3)$.

\subsubsection{NSGA-II}

\begin{defbox}[NSGA-II (Deb et al., 2000)]
It repairs the shortcomings of NSGA. Three pillars:
\textbf{(1) fast non-dominated sorting} into fronts (complexity $O(kN^2)$),
\textbf{(2) the crowding distance} instead of the parameter
$d_{\mathrm{share}}$,
\textbf{(3) elitism} --- the parents and the offspring are merged and sorted,
and the better half forms the new population.
\end{defbox}

\begin{defbox}[Crowding distance]
For each front separately: for each criterion $m$ sort the individuals according
to $f_m$, assign $\infty$ to the extreme individuals and add to the others the
normalised distance of their neighbours:
\[
d_i \;=\; \sum_{m=1}^{k}
\frac{f_m(i+1) - f_m(i-1)}{f_m^{\max} - f_m^{\min}} .
\]
It expresses the \uv{perimeter of the box} formed by the nearest neighbours ---
a large value means a sparsely occupied region, which we want to preserve.
\end{defbox}

\begin{defbox}[Crowded comparison operator]
An individual $i$ is better than $j$ if and only if
(a) $i$ lies in a \emph{lower} front ($rank_i < rank_j$), or
(b) the fronts are the same and $i$ has a \emph{larger} crowding distance.
No scalar fitness is therefore computed at all --- only these two numbers are
compared, typically in a binary tournament.
\end{defbox}

\begin{algorithm}[!ht]
\caption{NSGA-II (one generation)}
\begin{algorithmic}[1]
\State $R_t \gets P_t \cup Q_t$ \Comment{parents and offspring, $2N$ individuals in total}
\State $(F_1, F_2, \dots) \gets$ non-dominated sorting of $R_t$
\State $P_{t+1} \gets \emptyset$; $i \gets 1$
\While{$|P_{t+1}| + |F_i| \le N$}
  \State compute the crowding distance in $F_i$; $P_{t+1} \gets P_{t+1} \cup F_i$; $i \gets i+1$
\EndWhile
\State sort $F_i$ in decreasing order of the crowding distance and fill $P_{t+1}$ up to $N$ individuals
\State $Q_{t+1} \gets$ tournament selection (crowded comparison) $+$ crossover (SBX) $+$ mutation
\end{algorithmic}
\end{algorithm}

\paragraph{A numerical example of the crowding distance.} The front
$A(1,10)$, $B(2,7)$, $C(4,5)$, $D(8,2)$; the ranges: $f_1 \in [1,8]$
(a span of 7), $f_2 \in [2,10]$ (a span of 8). The extreme points $A$ and $D$
receive $\infty$ (they are extreme in both criteria).
\[
d_B = \frac{f_1(C)-f_1(A)}{7} + \frac{f_2(A)-f_2(C)}{8}
    = \frac{4-1}{7} + \frac{10-5}{8} = 0.43 + 0.63 = 1.05,
\]
\[
d_C = \frac{f_1(D)-f_1(B)}{7} + \frac{f_2(B)-f_2(D)}{8}
    = \frac{8-2}{7} + \frac{7-2}{8} = 0.86 + 0.63 = 1.48 .
\]
If one of $B$, $C$ has to be chosen, $C$ wins --- it lies in a more sparsely
occupied part of the front.

\paragraph{NSGA-III.} For problems with many criteria ($k \ge 4$,
\emph{many-objective}), where the crowding distance fails (almost everything is
non-dominated). The crowding distance is replaced by a set of evenly distributed
\textbf{reference points} (their number roughly corresponds to the population
size, $N = \binom{p+k-1}{p}$ for $p$ divisions). After the non-dominated sorting
those reference directions that have the fewest solutions assigned to them are
filled preferentially; if a direction has no solution, the individual with the
smallest perpendicular distance from the corresponding reference vector
survives.

\subsubsection{Other important algorithms}

\noindent\mimoq{the slides stop at NSGA-II}

\textbf{SPEA2} uses an external \textbf{archive} of non-dominated solutions; the fitness of
an individual is the sum of the \emph{strengths} of all the individuals that dominate it
(strength = the number of individuals the given individual dominates), plus a density derived
from the distance to the $\sigma$-th nearest neighbour.
\textbf{MOEA/D} decomposes the problem into $\mu$ \emph{scalar} subproblems with evenly
distributed weight vectors (Chebyshev scalarisation); every individual corresponds to one
subproblem and cooperates only with its neighbourhood --- hence the low complexity and good
scalability. \textbf{SMS-EMOA} selects directly by the contribution to the
\emph{hypervolume} (the best-founded theoretically, but exponentially expensive in the number
of criteria). \textbf{NSGA-III} replaces the crowding distance by evenly distributed
\textbf{reference points} and targets many-objective problems ($k \ge 4$), where the crowding
distance fails.

\subsubsection{Metrics of approximation quality}

\noindent\mimo{}

\begin{defbox}[Hypervolume (the $S$-metric)]
The volume of the region dominated by the front found and bounded by a \emph{reference point}
$r$ (chosen so that it is dominated by all the solutions). The only common metric that is
\emph{strictly monotone with respect to dominance} and that does not require knowledge of the
true Pareto front.

\emph{Example:} the front $\{(2,7), (4,5), (8,2)\}$, $r = (10,10)$. By vertical strips:
$(8,2)$ gives $2\cdot8 = 16$, $(4,5)$ gives $4\cdot5 = 20$, $(2,7)$ gives $2\cdot3 = 6$;
in total $HV = 42$. Adding a dominated point does not change the value.
\end{defbox}

Furthermore \emph{GD} (the average distance of the points found from the true front --- it
measures convergence), \emph{IGD} (the other way round; it measures both convergence and
coverage) and \emph{spread} (the evenness of the distribution); both GD and IGD require
knowledge of the true front.

\subsubsection{An overall comparison of MOEAs}

\begin{center}
\small
\begin{tabular}{@{}llll@{}}
\toprule
\textbf{algorithm} & \textbf{selection principle} & \textbf{diversity} & \textbf{note} \\
\midrule
VEGA & subpopulations by $f_i$ & none & converges to the extremes \\
NSGA & fronts + dummy fitness & fitness sharing ($d_{\mathrm{share}}$) & no elitism, $O(kN^3)$ \\
NSGA-II & fronts + crowding & crowding distance & elitism, $O(kN^2)$, standard \\
NSGA-III & fronts + reference points & reference directions & many-objective \\
SPEA2 & strength of dominators & $\sigma$-th neighbour, truncation & external archive \\
MOEA/D & scalarised subproblems & even weight vectors & neighbourhoods, scales \\
SMS-EMOA & hypervolume contribution & implicit & best-founded, expensive \\
\bottomrule
\end{tabular}
\end{center}

\subsection{Exam summary}

\begin{itemize}
\item \textbf{Michigan}: an individual = a single rule, the solution is the
  whole population, the reward has to be distributed (credit assignment).
  \textbf{Pittsburgh}: an individual = a whole set of rules, the fitness is
  straightforward, the operators are complex (GIL).
\item Holland's LCS: conditions over $\{0,1,*\}$, internal messages and
  receptors, the bucket brigade (a rule pays its predecessor for the right to
  act).
\item ZCS: no messages; the match set $M$ $\to$ roulette wheel selection of the
  action $\to$ the action set $A$; reinforcement of $A$ and $A_{-1}$, a tax for
  $M \setminus A$; the cover operator.
\item XCS: a classifier $(c,a,p,\epsilon,F)$, \textbf{fitness based on the
  accuracy of the prediction} $\kappa = \alpha(\epsilon/\epsilon_0)^{-\nu}$,
  updates by Q-learning, a niche GA in the action set; it evolves a complete and
  minimal input--output map. Variants: UCS (supervised learning), XCSF (function
  approximation).
\item Be able to explain why accuracy as the fitness solves the shortcomings of
  \uv{strength-based} systems.
\item What is evolved: the weights / the topology / both / the hyperparameters
  of the learning. The advantages of EAs: no gradient, recurrent networks, RL,
  parallelisation.
\item The problem of competing conventions (the permutation problem) --- without
  precautions, crossing networks is destructive.
\item Encodings of the architecture: direct (a matrix), grammar-based (Kitano),
  cellular (Gruau, a GP program for growing the network); GNARL = EP over
  graphs, structural as well as parametric mutations controlled by a
  temperature.
\item NEAT: a list of edges with \textbf{innovation numbers} (only matching
  genes are crossed), \textbf{speciation} according to the similarity of the
  genomes with shared fitness (protection of new topologies),
  \textbf{complexification} from a minimal network.
\item HyperNEAT: a substrate $S:\mathbb{R}^4 \to \mathbb{R}$ generating the
  weights from the coordinates of the neurons, represented by a CPPN (a DAG with
  various activation functions) evolved by NEAT; it captures symmetries and
  regularities.
\item NAS: the space (linear / a DAG of cells / a supernetwork), the strategy
  (EA, RL, Bayesian, gradient-based), the performance estimate (weight sharing
  --- ENAS); DeepNEAT and CoDeepNEAT (modules + blueprints); ES for deep RL.
\item The knapsack: a bit map, trivial operators, a problem with the capacity
  constraint. Three strategies for constraints: \textbf{a penalty} (the choice
  of the weight $\alpha$), \textbf{a repair} (Lamarckian/Baldwinian),
  \textbf{a decoder} (\uv{add it if it fits} --- every genotype is feasible);
  a fourth possibility = a multi-objective formulation.
\item The TSP: the fitness is trivial, the problem is the encoding. Do not use
  the adjacency and the ordinal representation; use the permutation one +
  PMX/OX/CX/ER, or possibly the matrix one. Initialisation: nearest neighbour,
  edge insertion. Mutation: inversion (2-opt).
\item SAT: the fitness = the number of satisfied clauses ($f$ alone is not
  enough), weighting of the clauses, a refining function; the representations
  are bit-based, real-valued, clause-based, path-based; WalkSAT as a local
  heuristic.
  With the real-valued representation ($x \mapsto (1-y)^2$,
  $\neg x \mapsto (1+y)^2$, minimisation) beware: the values are
  \emph{penalties}, so a disjunction (a clause) is a \textbf{product} and a
  conjunction (the formula) a \textbf{sum}.
\item Scheduling: a direct matrix encoding, hard constraints to be handled in
  the operators, soft ones in the fitness; job shop --- a permutation with a
  decoder vs.\ a direct encoding.
\item The general rule: in combinatorial problems what matters is the encoding
  and operators that respect the structure of the problem; hybridisation with a
  local heuristic (2-opt, WSAT) is almost always advantageous.
\item Dominance ($x \prec y$: not worse in all the criteria and better in at
  least one), the Pareto set and the Pareto front; the twofold goal =
  convergence + diversity.
\item Scalarisation: linear (simple, but it will not find the non-convex parts
  of the front), $\varepsilon$-constraint, Chebyshev.
\item NSGA-II: non-dominated sorting into fronts, the crowding distance
  $d_i = \sum_m (f_m(i+1)-f_m(i-1))/(f_m^{\max}-f_m^{\min})$, the extreme points
  $\infty$; the crowded comparison (a lower front, and in case of a tie a larger
  crowding distance); elitism through the merge $P_t \cup Q_t$. Be able to
  compute the crowding distance on an example.
\item SPEA2 (an archive, the strength of the dominators, the density via the
  $\sigma$-th neighbour), MOEA/D (Chebyshev scalarisation, weight vectors,
  neighbourhoods), SMS-EMOA (the hypervolume contribution), NSGA-III (reference
  points for many-objective problems).
\item Hypervolume: the volume dominated by the front with respect to a reference
  point; the only common metric that is monotone with respect to dominance and
  that does not need the true front.
  Furthermore GD, IGD, spread, the $\varepsilon$-indicator.
\end{itemize}

%=====================================================================
\section*{Overall comparison and typical examination questions}
\addcontentsline{toc}{section}{Overall comparison and typical examination questions}
\label{sec:zaver}
%=====================================================================
\subsection*{A map of the field}

\begin{itemize}
\item \textbf{Evolutionary algorithms} (a population, fitness, crossover,
  mutation, selection): genetic algorithms (binary, integer, permutation,
  real-valued), evolution strategies, differential evolution, evolutionary
  programming, genetic programming (tree-based, linear, Cartesian, grammatical
  evolution), learning classifier systems, estimation of distribution algorithms
  (EDA).
\item \textbf{Nature-inspired methods without evolution}: swarm algorithms
  (PSO, ACO, ABC, firefly).
\item \textbf{Metaheuristics not inspired by biology}: tabu search,
  scatter search, novelty search, simulated annealing, memetic algorithms.
\item \textbf{Application areas} with their own representations and operators:
  combinatorial optimisation, neuroevolution, reinforcement learning,
  multi-objective optimisation, AutoML/NAS.
\end{itemize}

\subsection*{An overall comparison table of the algorithms}

\begin{center}
\small
\begin{tabular}{@{}p{2.5cm}p{2.9cm}p{3.1cm}p{3.1cm}p{3.1cm}@{}}
\toprule
\textbf{algorithm} & \textbf{representation} & \textbf{variation} &
\textbf{selection} & \textbf{what to remember} \\
\midrule
SGA & binary string & 1-point crossover, bit-flip & roulette wheel,
  generational & schema theory, BBH \\
ES & $\mathbb{R}^n$ + $\sigma$ & Gaussian mutation, recombination &
  random parents, $(\mu,\lambda)$/$(\mu{+}\lambda)$ & the 1/5 rule,
  self-adaptation \\
CMA-ES & $\mathbb{R}^n$, the model $\mathcal N(m,\sigma^2 C)$ & sampling
  from the distribution & rank-based, the $\mu$ best & evolution paths,
  rank-1/rank-$\mu$ \\
DE & $\mathbb{R}^n$ & $\vec x_a + F(\vec x_b - \vec x_c)$, binomial crossover &
  parent--offspring duel & the scale from the population \\
EP & domain-specific (an automaton) & mutation only & tournament of parents
  and offspring & genotype = phenotype \\
GP & syntax tree & subtree exchange, mutation & tournament &
  bloat, ADFs \\
PSO & $\mathbb{R}^n$ + velocity & the equations of motion & none (memory only) &
  $w$, $c_1$, $c_2$, gbest/lbest \\
ACO & construction of a path & pheromone + heuristic & implicit &
  $\tau^\alpha\eta^\beta$, evaporation \\
SA & a single solution & a neighbour & Metropolis & the cooling schedule \\
MA & arbitrary & EA operators + local search & arbitrary &
  Lamarck vs.\ Baldwin \\
NSGA-II & arbitrary & SBX + polynomial mutation & fronts + crowding &
  elitism, the crowding distance \\
NEAT & list of edges & structural mutations, crossover by ID & shared fitness
  within a species & innovation numbers, speciation \\
\bottomrule
\end{tabular}
\end{center}

\subsection*{Cross-cutting topics --- what is asked across the questions}

\begin{itemize}
\item \textbf{Exploration vs.\ exploitation} --- how each algorithm handles it:
  a GA (crossover vs.\ the selection pressure), an ES ($\sigma$ and the 1/5
  rule), SA (the temperature), PSO ($w$ and the topology), ACO ($q_0$,
  evaporation), tabu search (the tabu list), novelty search (exploration only).
\item \textbf{Maintaining diversity} --- fitness sharing, crowding, speciation,
  the island model, non-random mating, restarts, hypermutation.
\item \textbf{Parameter adaptation} --- tuning vs.\ control; fixed, adaptive and
  self-adaptive strategies; where each of them occurs in which algorithm.
\item \textbf{Constraint handling} --- penalties, repair, decoders, closed
  operators, a multi-objective formulation.
\item \textbf{Lamarckism vs.\ Baldwinism} --- memetic algorithms, the repair of
  individuals, neuroevolution with the retraining of the weights.
\item \textbf{The theoretical background} --- the schema theorem and its
  criticism, Vose's model, Markov chains and the convergence of an elitist GA,
  No Free Lunch, the convergence of SA.
\end{itemize}

\subsection*{Typical examination questions}

\begin{enumerate}
\item Describe the general scheme of an evolutionary algorithm and list the
  types of selection together with their properties. Compute roulette wheel
  selection for a given population.
\item State and prove the schema theorem. What are the order and the defining
  length? What are the weaknesses of the theorem?
\item Explain the building blocks hypothesis and implicit parallelism.
  What is a deceptive problem?
\item Describe Vose's model of the SGA: what are the vectors $p$ and $s$, what
  does the operator $\mathcal{F}$ do, what does $\mathcal{M}$ do, and what fixed
  points do they have? What does the finite-population model as a Markov chain
  look like?
\item What are estimation of distribution algorithms (EDA)? What is their
  relation to Vose's view of a GA and how do UMDA/PBIL differ from BOA?
\item What is the difference between a $(\mu,\lambda)$- and a
  $(\mu+\lambda)$-ES? Explain the 1/5 rule and the self-adaptation of $\sigma$.
  How does CMA-ES work?
\item Describe one step of differential evolution including the formulas and a
  numerical example. What are the mutation variants?
\item How does tree-based genetic programming work? What is bloat and how is it
  fought? What are ADFs for?
\item Explain grammatical evolution and Cartesian GP --- what is the genotype
  and what is the phenotype?
\item What is coevolution, what types does it have and what pathologies? How is
  fitness evaluated?
\item What is open-ended evolution, what are its hallmarks and why do artificial
  systems \uv{get stuck}? How does novelty search work and why does it pay off
  to throw the goal away in deceptive problems?
\item Explain the iterated prisoner's dilemma, the TFT strategy and the
  properties of successful strategies.
\item Write down the PSO equations and explain the meaning of the individual
  terms. What is the difference between the gbest and the lbest topology?
\item Describe ACO: the transition rule, the pheromone update, the difference
  between AS and ACS.
\item Explain simulated annealing, the Metropolis criterion and the cooling
  schedule. When does it converge?
\item What is a memetic algorithm? Explain the difference between Lamarckian and
  Baldwinian learning.
\item Compare the Michigan and the Pittsburgh approach. How does XCS work and
  why is its fitness based on accuracy?
\item Describe NEAT: what are the innovation numbers and the speciation for?
  What does HyperNEAT add?
\item How is the TSP solved by an evolutionary algorithm? Describe PMX, OX, CX
  and ER.
\item How are constraints handled in an EA (using the knapsack problem as an
  example)?
\item Define dominance and the Pareto front. Describe NSGA-II and compute the
  crowding distance on an example. What is the hypervolume?
\end{enumerate}

\subsection*{Concluding summary}

\begin{itemize}
\item Evolutionary and swarm algorithms are \textbf{robust metaheuristics} for
  problems where exact and gradient methods fail: black-box, non-differentiable,
  discrete, structured, multi-objective and dynamic problems.
\item Their strength does not lie in one \uv{best} algorithm (No Free Lunch),
  but in \textbf{the ability to adapt the representation and the operators to
  the given domain} and in the possibility of \textbf{hybridisation} with local
  heuristics and domain knowledge.
\item The theory (schemata, Vose's model, convergence theorems, runtime
  analysis) provides intuition and asymptotic guarantees, but the practical
  design remains to a large extent an experimental discipline.
\end{itemize}

\end{document}
